%% 3-2 場所を数で表す：座標（第12回：組 V3-1）  % noruby
%% ①数直線を2本に ②原点・x軸・y軸・(x, y) ③マイナスの座標と4つの区画 ④＋の向きと原点は自分で決め、決めたら変えない  % noruby
%% ⑤点を動かす ⑥下つき文字 x_1（かけ算ではない） ⑦絶対値と、横・たてに並んだ2点の距離（ななめは 4-4）  % noruby
\makeatletter
%% ---- この章だけで使う絵の部品（頭文字 \bkthrc…） ----
% 座標の方眼  \bkthrcgrid{xmin}{xmax}{ymin}{ymax}{1目もりのmm}{目もりの数字 1:全部 2:偶数だけ 0:なし}  原点が (0,0)  % noruby
\newcommand\bkthrcgrid[6]{%
  \begin{scope}[x=#5mm, y=#5mm]
  \draw[bkLine, line width=0.4pt] (#1,#3) grid (#2,#4);
  \draw[-{Stealth[length=4pt]}, bkInk, line width=0.8pt] (#1-0.4,0) -- (#2+0.7,0);
  \draw[-{Stealth[length=4pt]}, bkInk, line width=0.8pt] (0,#3-0.4) -- (0,#4+0.7);
  \node[anchor=west, font=\small, inner sep=1pt] at (#2+0.75,0) {$x$};
  \node[anchor=south, font=\small, inner sep=1pt] at (0,#4+0.75) {$y$};
  \ifnum#6=1
    \foreach \k in {#1,...,#2} {\ifnum\k=0\else\node[anchor=north, font=\footnotesize, fill=white, inner sep=0.6pt] at (\k,-0.15) {$\k$};\fi}
    \foreach \k in {#3,...,#4} {\ifnum\k=0\else\node[anchor=east, font=\footnotesize, fill=white, inner sep=0.6pt] at (-0.15,\k) {$\k$};\fi}
  \fi
  \ifnum#6=2
    \foreach \k in {#1,...,#2} {\ifodd\k\else\ifnum\k=0\else\node[anchor=north, font=\footnotesize, fill=white, inner sep=0.6pt] at (\k,-0.15) {$\k$};\fi\fi}
    \foreach \k in {#3,...,#4} {\ifodd\k\else\ifnum\k=0\else\node[anchor=east, font=\footnotesize, fill=white, inner sep=0.6pt] at (-0.15,\k) {$\k$};\fi\fi}
  \fi
  \node[anchor=north east, font=\footnotesize, fill=white, inner sep=0.6pt] at (-0.12,-0.12) {O};
  \end{scope}}
% 点  \bkthrcpt{(x,y)}{1目もりのmm}{名前}{名前の向き}{色}  % noruby
\newcommand\bkthrcpt[5]{%
  \begin{scope}[x=#2mm, y=#2mm]
  \fill[#5] #1 circle (1.2mm);
  \node[anchor=#4, font=\sffamily\bfseries\small, text=#5!80!black, inner sep=1.6mm] at #1 {#3};
  \end{scope}}
% 教室の机  \bkthrcdesk{(x,y)}{色}  はば 9 mm・おく 6 mm  % noruby
\newcommand\bkthrcdesk[2]{\filldraw[fill=#2, draw=bkGray!80!black, line width=0.5pt, rounded corners=0.6mm] #1 rectangle ++(9,6);}
% 宝箱  \bkthrctreasure{(x,y)}  % noruby
\newcommand\bkthrctreasure[1]{\begin{scope}[shift={#1}]
  \filldraw[fill=bkOrange!70, draw=bkOrange!50!black, line width=0.5pt] (-2.6,-1.8) rectangle (2.6,1.2);
  \filldraw[fill=bkOrange!85, draw=bkOrange!50!black, line width=0.5pt] (-2.6,1.2) arc (180:0:2.6 and 1.3) -- cycle;
  \fill[bkYellow] (-0.6,0.2) rectangle (0.6,1.6);
\end{scope}}
% 本文の中の小見出し（それだけで1行）  % noruby
\newcommand\bkthrchead[1]{\par\addvspace{4pt}\Needspace{4\baselineskip}\noindent\textbf{#1}\par\nobreak}
\makeatother

\bkchapter{3}{2}{\R{場所}{ば|しょ}を\R{数}{かず}で\R{表}{あらわ}す：\R{座標}{ざ|ひょう}}{\bklevelrange{1}{2}}{\R{場所}{ば|しょ}を\R{数}{かず}で\R{表}{あらわ}す：\R{座標}{ざ|ひょう}}
\bkmichishirube{1-6,2-2,2-15}{3-3,3-4,4-4,4-10,4-13,4-15}

\begin{goal}
\begin{itemize}[label={\color{\bkTT}●}, leftmargin=1.6\zw, itemsep=0pt, topsep=1pt]
\item \bklevel{1}\enspace \R{原点}{げん|てん}と2\R{本}{ほん}の\R{軸}{じく}を\R{決}{き}めて、\R{点}{てん} $(x,\ y)$ をかいたり、\R{点}{てん}の\R{座標}{ざ|ひょう}を\R{読}{よ}んだりできる。
\item \bklevel{2}\enspace マイナスの\R{座標}{ざ|ひょう}と4つの\R{区画}{く|かく}、\R{点}{てん}を\R{動}{うご}かしたときの\R{座標}{ざ|ひょう}がわかる。
\item \bklevel{2}\enspace \R{絶対値}{ぜっ|たい|ち}を\R{使}{つか}って、\R{横}{よこ}やたてに\R{並}{なら}んだ2\R{点}{てん}の\R{距離}{きょ|り}を\R{求}{もと}められる。
\end{itemize}
\end{goal}

\begin{hatena}
\noindent\begin{minipage}[c]{0.5\linewidth}
\R{今日}{きょう}、\R{教室}{きょう|しつ}の\R{席}{せき}がえがありました。\R{休}{やす}んだ\R{友}{とも}だちに、\R{電話}{でん|わ}で\R{新}{あたら}しい\R{席}{せき}を\R{伝}{つた}えます。\par\vspace{2pt}
「まどから4\R{列}{れつ}め」だけで、\R{席}{せき}は1つに\R{決}{き}まるかな？\quad 「どこ」を、いくつの\R{数}{かず}で\R{伝}{つた}えればいいかな？
\end{minipage}\hfill
\begin{minipage}[c]{0.47\linewidth}\centering
\begin{tikzpicture}[x=1mm, y=1mm]
  % 黒板（前）とまど（左）
  \fill[bkGreen!55!black, rounded corners=0.6mm] (8,46) rectangle (51,49.5);
  \node[anchor=south, font=\sffamily\footnotesize, inner sep=0.8pt] at (29.5,50) {\R{黒板}{こく|ばん}（\R{前}{まえ}）};
  \draw[bkBlue!70, line width=1.6pt] (0,2) -- (0,40);
  \node[anchor=south, font=\sffamily\footnotesize, rotate=90, inner sep=1pt] at (-1.5,21) {まど};
  \foreach \c in {0,...,5} {\foreach \r in {0,...,4} {\filldraw[fill=white, draw=bkGray!80!black, line width=0.5pt, rounded corners=0.6mm] ({4+\c*8.6},{36-\r*8.4}) rectangle ++(7.4,6);}}
  \filldraw[fill=bkOrange!60, draw=bkGray!80!black, line width=0.5pt, rounded corners=0.6mm] (29.8,19.2) rectangle ++(7.4,6);
  \node[font=\sffamily\bfseries\small, text=bkOrange!80!black] at (33.5,22.2) {?};
  \foreach \c in {0,...,5} {\pgfmathtruncatemacro\n{\c+1}\node[font=\sffamily\footnotesize, text=bkGray, anchor=north] at ({7.7+\c*8.6},-0.5) {\n};}
\end{tikzpicture}
\end{minipage}
\end{hatena}

\begin{omoidasou}
\bkrecall{1-6}{\R{数直線}{すう|ちょく|せん}は、0から\R{同}{おな}じ\R{間}{ま}かくで\R{数}{かず}を\R{並}{なら}べたもの。\R{右}{みぎ}へ\R{進}{すす}むほど\R{大}{おお}きい。}
\bkrecall{2-2}{\R{負}{ふ}の\R{数}{かず}は0の\R{左}{ひだり}に\R{並}{なら}ぶ。マイナスは「\R{反対}{はん|たい}の\R{向}{む}き」。}
\bkrecall{2-15}{グラフの\R{点}{てん}は、0から\R{右}{みぎ}へ2、そこから\R{上}{うえ}へ4\R{進}{すす}んだ\R{所}{ところ}にかく（$x=2$、$y=4$）。}
\bkrecall{2-3}{\R{変}{か}わった\R{量}{りょう}は「あと－まえ」。$\Delta x$ と\R{書}{か}く。}
\end{omoidasou}

\begin{yattemiyou}[\R{宝}{たから}さがしゲームをしよう]
\noindent\begin{minipage}[c]{0.54\linewidth}
\begin{enumerate}[label={\bkcircled[bkGreen]{\arabic*}}, leftmargin=2.4\zw, itemsep=1pt, topsep=0pt]
\item \R{方眼紙}{ほう|がん|し}のまん\R{中}{なか}に\R{木}{き}を1\R{本}{ぽん}かき、\R{上}{うえ}を\R{北}{きた}とする。\R{友}{とも}だちに\R{見}{み}せないように、\R{宝}{たから}の\R{場所}{ば|しょ}に\R{印}{しるし}をつけよう。
\item \R{友}{とも}だちに、ことばだけで\R{宝}{たから}の\R{場所}{ば|しょ}を\R{伝}{つた}えよう。\R{友}{とも}だちは\R{自分}{じ|ぶん}の\R{方眼紙}{ほう|がん|し}に\R{印}{しるし}をつける。
\item 2\R{人}{にん}の\R{紙}{かみ}を\R{重}{かさ}ねて、\R{印}{しるし}が\R{同}{おな}じ\R{場所}{ば|しょ}にあるか\R{確}{たし}かめよう。\R{伝}{つた}えるのに\R{何}{なに}が\R{必要}{ひつ|よう}だったかな？
\end{enumerate}
\end{minipage}\hfill
\begin{minipage}[c]{0.42\linewidth}\centering
\begin{tikzpicture}[x=1mm, y=1mm]
  \draw[bkLine, line width=0.4pt] (0,0) grid[step=5.5] (44,33);
  \bktree{(22,16.5)}{8.5}
  \bkthrctreasure{(38.5,27.5)}
  \draw[bkRed, line width=0.9pt, densely dashed, -{Stealth[length=3pt]}] (23.5,16.5) -- (38.5,16.5);
  \draw[bkRed, line width=0.9pt, densely dashed, -{Stealth[length=3pt]}] (38.5,16.5) -- (38.5,25);
  \node[anchor=north, font=\sffamily\footnotesize, text=bkRed!80!black, fill=white, inner sep=0.6pt] at (30.25,15.3) {\R{東}{ひがし}へ3};
  \node[anchor=west, font=\sffamily\footnotesize, text=bkRed!80!black, fill=white, inner sep=0.6pt] at (39.5,21) {\R{北}{きた}へ2};
  \node[anchor=south, font=\sffamily\footnotesize, inner sep=0.6pt] at (6,34) {↑\R{北}{きた}};
\end{tikzpicture}
\end{minipage}
\end{yattemiyou}

\noindent 「\R{宝}{たから}は\R{東}{ひがし}のほう」だけでは、\R{伝}{つた}わりません。うまく\R{伝}{つた}えるには、\R{次}{つぎ}の3つが\R{必要}{ひつ|よう}だったはずです。

\begin{enumerate}[label={\bkcircled{\arabic*}}, leftmargin=2.4\zw, itemsep=1pt, topsep=2pt]
\item どこから\R{数}{かぞ}えるか（\R{木}{き}の\R{場所}{ば|しょ}）
\item どちら\R{向}{む}きに\R{数}{かぞ}えるか（\R{東}{ひがし}と\R{北}{きた}）
\item いくつ\R{進}{すす}むか（\R{東}{ひがし}へ3マス、\R{北}{きた}へ2マス）
\end{enumerate}

\noindent しかも、③の\R{数}{かず}は「\R{東}{ひがし}へ3」と「\R{北}{きた}へ2」の2つがいりました。この\R{章}{しょう}では、\R{平}{たい}らな\R{紙}{かみ}や\R{画面}{が|めん}の\R{上}{うえ}の\R{場所}{ば|しょ}を、2つの\R{数}{かず}で\R{表}{あらわ}す\R{方法}{ほう|ほう}を\R{学}{まな}びます。

%% ---------------- 数学の道具 ----------------
\bksection{doguu}{\R{数学}{すう|がく}の\R{道具}{どう|ぐ}}

\bksubsection{1}{\R{数直線}{すう|ちょく|せん}を、たてと\R{横}{よこ}の2\R{本}{ほん}にする}

\label{bk:3-2:jiku}\bkthrchead{1\R{本}{ぽん}の\R{線}{せん}の\R{上}{うえ}なら、\R{数}{かず}は1つでよい}
まっすぐな\R{線路}{せん|ろ}を\R{走}{はし}る\R{電車}{でん|しゃ}の\R{場所}{ば|しょ}は、「\R{駅}{えき}から\R{東}{ひがし}へ\qty{3}{km}」のように1つの\R{数}{かず}で\R{言}{い}えます。\R{駅}{えき}を0、\R{東}{ひがし}を＋の\R{向}{む}きにした\R{数直線}{すう|ちょく|せん}（2-2）を\R{線路}{せん|ろ}にそって\R{置}{お}けば、\R{東}{ひがし}へ\qty{3}{km}は3、\R{西}{にし}へ\qty{2}{km}は $-2$ です。

\begin{bkfigure}
\begin{tikzpicture}[x=1mm, y=1mm]
  \draw[bkGray!70, line width=2.4pt] (-60,0) -- (60,0);
  \foreach \k in {-5,...,5} {\draw[bkInk, line width=0.5pt] (\k*11,-1.4) -- (\k*11,1.4);
    \node[anchor=north, font=\footnotesize, inner sep=0pt] at (\k*11,-2.6) {$\k$};}
  \draw[-{Stealth[length=4pt]}, bkInk, line width=0.8pt] (57,0) -- (64,0);
  \node[anchor=west, font=\sffamily\footnotesize] at (64.5,0) {km};
  \fill[bkGreen!60!black] (-2,1.5) rectangle (2,6);
  \node[anchor=south, font=\sffamily\footnotesize] at (0,6.5) {\R{駅}{えき}};
  \bkcar{(28.5,1.5)}{0.3}{bkBlue}
  \bkcar{(-27.5,1.5)}{0.3}{bkOrange}
  \node[anchor=south, font=\sffamily\footnotesize, text=bkBlue!80!black] at (33,7.5) {3};
  \node[anchor=south, font=\sffamily\footnotesize, text=bkOrange!80!black] at (-22,7.5) {$-2$};
  \node[anchor=north, font=\sffamily\footnotesize, text=bkGray] at (-55,-7) {\R{西}{にし}};
  \node[anchor=north, font=\sffamily\footnotesize, text=bkGray] at (55,-7) {\R{東}{ひがし}};
\end{tikzpicture}
\bkfigcaption{\R{線路}{せん|ろ}の\R{上}{うえ}の\R{場所}{ば|しょ}は、1つの\R{数}{かず}で\R{表}{あらわ}せる。}
\end{bkfigure}

\bkthrchead{\R{平}{たい}らな\R{所}{ところ}では、\R{数}{かず}が2ついる}
\noindent\begin{minipage}[c]{0.54\linewidth}
\R{町}{まち}の\R{中}{なか}のように、\R{東西}{とう|ざい}にも\R{南北}{なん|ぼく}にも\R{動}{うご}ける\R{所}{ところ}ではどうでしょう。「\R{学校}{がっ|こう}から\R{東}{ひがし}へ3」と\R{言}{い}われても、\R{右}{みぎ}の\R{図}{ず}の\R{点線}{てん|せん}の\R{上}{うえ}のどこなのか、わかりません。\R{北}{きた}へ\R{何}{なん}マス\R{進}{すす}むかも\R{言}{い}わないと、\R{場所}{ば|しょ}が1つに\R{決}{き}まらないのです。
\end{minipage}\hfill
\begin{minipage}[c]{0.42\linewidth}\centering
\begin{tikzpicture}[x=1mm, y=1mm]
  \draw[bkLine, line width=0.4pt] (0,0) grid[step=7] (49,35);
  \fill[bkNavy] (12.5,12.5) rectangle (15.5,15.5);
  \node[anchor=north east, font=\sffamily\footnotesize, fill=white, inner sep=0.6pt] at (12.3,12.3) {\R{学校}{がっ|こう}};
  \draw[bkOrange, line width=1.4pt, dash pattern=on 2.5pt off 1.5pt] (35,0) -- (35,35);
  \foreach \yy in {0,7,...,35} \fill[bkOrange] (35,\yy) circle (1);
  \node[anchor=south, font=\sffamily\footnotesize, text=bkOrange!80!black, fill=white, inner sep=0.6pt] at (35,36) {\R{東}{ひがし}へ3};
\end{tikzpicture}
\end{minipage}

\par\vspace{2mm}
そこで、\R{数直線}{すう|ちょく|せん}を2\R{本}{ほん}\R{使}{つか}います。\R{横}{よこ}の\R{数直線}{すう|ちょく|せん}と、たての\R{数直線}{すう|ちょく|せん}を、どちらも0の\R{所}{ところ}で\R{直角}{ちょっ|かく}に\R{交}{まじ}わるように\R{置}{お}きます。\R{横}{よこ}の\R{数直線}{すう|ちょく|せん}は\R{右}{みぎ}が＋、たての\R{数直線}{すう|ちょく|せん}は\R{上}{うえ}が＋です。

\noindent\begin{minipage}[c]{0.50\linewidth}\centering
\begin{tikzpicture}[x=1mm, y=1mm]
  \bkthrcgrid{-4}{4}{-3}{3}{7}{1}
  \node[anchor=south west, font=\sffamily\footnotesize, text=bkBlue!80!black, fill=white, inner sep=0.8pt] at (14.5,1.2) {$x$ \R{軸}{じく}};
  \node[anchor=south west, font=\sffamily\footnotesize, text=bkBlue!80!black, fill=white, inner sep=0.8pt] at (1.2,15.5) {$y$ \R{軸}{じく}};
  \node[anchor=north west, font=\sffamily\footnotesize, text=bkRed!80!black, fill=white, inner sep=0.8pt] at (2,-9) {\R{原点}{げん|てん} O};
  \draw[bkRed, line width=0.6pt] (0.8,-0.8) -- (3,-8.6);
\end{tikzpicture}
\end{minipage}\hfill
\begin{minipage}[c]{0.46\linewidth}
\R{横}{よこ}の\R{数直線}{すう|ちょく|せん}を\textbf{$x$ \R{軸}{じく}}、たての\R{数直線}{すう|ちょく|せん}を\textbf{$y$ \R{軸}{じく}}、2\R{本}{ほん}をまとめて\textbf{\R{座標軸}{ざ|ひょう|じく}}といいます。2\R{本}{ほん}の\R{軸}{じく}が\R{交}{まじ}わる0の\R{点}{てん}を\textbf{\R{原点}{げん|てん}}といい、O と\R{書}{か}きます。O は、「はじまり」という\R{意味}{い|み}の\R{英語}{えい|ご} origin の\R{頭}{あたま}の\R{文字}{も|じ}で、\R{数字}{すう|じ}の0ではありません。\bkidx{げんてん}{\R{原点}{げん|てん}}
\end{minipage}

\bksubsection{2}{\R{点}{てん}の\R{場所}{ば|しょ}を2つの\R{数}{かず}で：\R{座標}{ざ|ひょう} $(x,\ y)$}

\noindent\label{bk:3-2:zahyo}\begin{minipage}[c]{0.44\linewidth}
\R{右}{みぎ}の\R{図}{ず}の\R{点}{てん} A は、\R{原点}{げん|てん} O から\R{右}{みぎ}へ3、そこから\R{上}{うえ}へ2\R{進}{すす}んだ\R{所}{ところ}にあります。2-15 のグラフで、\R{点}{てん}を「0から\R{右}{みぎ}へ2、そこから\R{上}{うえ}へ4」とかいたのと\R{同}{おな}じです。

「\R{右}{みぎ}へ3」は $x$ \R{軸}{じく}の3、「\R{上}{うえ}へ2」は $y$ \R{軸}{じく}の2です。この2つの\R{数}{かず}を、かっこの\R{中}{なか}に $x$ が\R{先}{さき}、$y$ が\R{後}{あと}の\R{順}{じゅん}に\R{並}{なら}べて、
\[ (3,\ 2) \]
と\R{書}{か}きます。
\end{minipage}\hfill
\begin{minipage}[c]{0.52\linewidth}\centering
\begin{tikzpicture}[x=1mm, y=1mm]
  \bkthrcgrid{-1}{5}{-1}{4}{9}{1}
  \draw[bkBlue, line width=1.2pt, -{Stealth[length=5pt]}] (0.5,1.2) -- (26.5,1.2);
  \draw[bkBlue, line width=1.2pt, -{Stealth[length=5pt]}] (27,1.6) -- (27,16.5);
  \node[anchor=south, font=\sffamily\footnotesize, text=bkBlue!80!black, fill=white, inner sep=0.6pt] at (13.5,2.4) {\R{右}{みぎ}へ3};
  \node[anchor=west, font=\sffamily\footnotesize, text=bkBlue!80!black, fill=white, inner sep=0.6pt] at (28.3,8.5) {\R{上}{うえ}へ2};
  \bkthrcpt{(3,2)}{9}{A$(3,\ 2)$}{south}{bkRed}
\end{tikzpicture}
\end{minipage}

\par\vspace{2mm}
\R{点}{てん}の\R{場所}{ば|しょ}を\R{表}{あらわ}すこの\R{数}{かず}の\R{組}{くみ}を、その\R{点}{てん}の\textbf{\R{座標}{ざ|ひょう}}といいます。\bkidx{ざひょう}{\R{座標}{ざ|ひょう}}3を\textbf{$x$ \R{座標}{ざ|ひょう}}、2を\textbf{$y$ \R{座標}{ざ|ひょう}}といいます。「\R{点}{てん} A の\R{座標}{ざ|ひょう}は $(3,\ 2)$」を、\R{短}{みじか}く A$(3,\ 2)$ とも\R{書}{か}きます。$(3,\ 2)$ は「さん、に」と\R{読}{よ}みます。

\begin{taisetsu}
\R{点}{てん} $(a,\ b)$ をかくには、\R{原点}{げん|てん} O から\R{右}{みぎ}へ $a$、そこから\R{上}{うえ}へ $b$ \R{進}{すす}む。\par
かっこの\R{中}{なか}は、いつも $x$ \R{座標}{ざ|ひょう}が\R{先}{さき}、$y$ \R{座標}{ざ|ひょう}が\R{後}{あと}。
\end{taisetsu}

\bkthrchead{\R{順番}{じゅん|ばん}を\R{入}{い}れかえると、ちがう\R{点}{てん}}
\noindent\begin{minipage}[c]{0.50\linewidth}\centering
\begin{tikzpicture}[x=1mm, y=1mm]
  \bkthrcgrid{-1}{4}{-1}{4}{9}{1}
  \bkthrcpt{(3,2)}{9}{$(3,\ 2)$}{west}{bkRed}
  \bkthrcpt{(2,3)}{9}{$(2,\ 3)$}{south}{bkBlue}
\end{tikzpicture}
\end{minipage}\hfill
\begin{minipage}[c]{0.46\linewidth}
$(3,\ 2)$ と $(2,\ 3)$ は、\R{使}{つか}っている\R{数}{かず}は\R{同}{おな}じでも、ちがう\R{場所}{ば|しょ}の\R{点}{てん}です。$(2,\ 3)$ は「\R{右}{みぎ}へ2、\R{上}{うえ}へ3」だからです。\R{座標}{ざ|ひょう}は、\R{数}{かず}の\R{組}{くみ}であるだけでなく、\R{並}{なら}べる\R{順番}{じゅん|ばん}にも\R{意味}{い|み}があります。
\end{minipage}

\bkthrchead{\R{点}{てん}から\R{座標}{ざ|ひょう}を\R{読}{よ}む}
\noindent\begin{minipage}[c]{0.44\linewidth}
\R{反対}{はん|たい}に、かいてある\R{点}{てん}の\R{座標}{ざ|ひょう}を\R{読}{よ}むには、\R{点}{てん}から\R{軸}{じく}へまっすぐな\R{線}{せん}をおろします。\R{点}{てん} B から\R{真下}{ま|した}へおろした\R{線}{せん}は $x$ \R{軸}{じく}の4に、\R{真横}{ま|よこ}へ\R{引}{ひ}いた\R{線}{せん}は $y$ \R{軸}{じく}の3に\R{当}{あ}たります。だから B$(4,\ 3)$ です。\R{進}{すす}む\R{道}{みち}すじ（\R{矢印}{や|じるし}）でも、\R{軸}{じく}へおろす\R{線}{せん}（\R{点線}{てん|せん}）でも、\R{同}{おな}じ2つの\R{数}{かず}が\R{見}{み}つかります。
\end{minipage}\hfill
\begin{minipage}[c]{0.52\linewidth}\centering
\begin{tikzpicture}[x=1mm, y=1mm]
  \bkthrcgrid{-1}{5}{-1}{4}{9}{1}
  \draw[bkRed, line width=0.9pt, dash pattern=on 2pt off 1.5pt] (36,27) -- (36,1.5);
  \draw[bkRed, line width=0.9pt, dash pattern=on 2pt off 1.5pt] (36,27) -- (1.5,27);
  \bkthrcpt{(4,3)}{9}{B}{south west}{bkGreen!80!black}
\end{tikzpicture}
\end{minipage}

\par\vspace{2mm}
\R{軸}{じく}の\R{上}{うえ}の\R{点}{てん}にも\R{座標}{ざ|ひょう}があります。$x$ \R{軸}{じく}の\R{上}{うえ}の4の\R{点}{てん}は、\R{上}{うえ}へは0\R{進}{すす}むので $(4,\ 0)$。$y$ \R{軸}{じく}の\R{上}{うえ}の3の\R{点}{てん}は $(0,\ 3)$。\R{原点}{げん|てん}は $(0,\ 0)$ です。

\R{点}{てん}が\R{目}{め}もりの\R{線}{せん}の\R{交}{まじ}わる\R{所}{ところ}にないこともあります。\R{右}{みぎ}へ2と3のまん\R{中}{なか}、\R{上}{うえ}へ1と2のまん\R{中}{なか}にある\R{点}{てん}なら、$(2.5,\ 1.5)$ です。\R{座標}{ざ|ひょう}は\R{整数}{せい|すう}だけでなく、\R{小数}{しょう|すう}や\R{分数}{ぶん|すう}にもなります。はかったデータの\R{点}{てん}（3-4）は、たいていこのような\R{点}{てん}です。

\begin{yattemiyou}[\R{点}{てん}をつないで、\R{絵}{え}をかこう]
\noindent\begin{minipage}[c]{0.6\linewidth}
\begin{enumerate}[label={\bkcircled[bkGreen]{\arabic*}}, leftmargin=2.4\zw, itemsep=1pt, topsep=0pt]
\item \R{方眼}{ほう|がん}のノートに、$x$ \R{軸}{じく}（$-4$〜4）と $y$ \R{軸}{じく}（$-5$〜6）をかき、\R{原点}{げん|てん} O を\R{書}{か}きこもう。
\item \R{次}{つぎ}の\R{点}{てん}を\R{順}{じゅん}にうって、\R{直線}{ちょく|せん}でつなごう。\R{最後}{さい|ご}の\R{点}{てん}は、はじめの\R{点}{てん}にもどす。\par
$(0,\ 5)$ → $(2,\ 2)$ → $(2,\ -2)$ → $(3,\ -4)$ → $(1,\ -3)$ → $(-1,\ -3)$ → $(-3,\ -4)$ → $(-2,\ -2)$ → $(-2,\ 2)$ → $(0,\ 5)$
\item \R{何}{なん}の\R{絵}{え}ができたかな？\quad \R{右}{みぎ}の\R{図}{ず}で\R{確}{たし}かめよう。
\item $(2,\ 2)$ と $(-2,\ 2)$ のように、$x$ \R{座標}{ざ|ひょう}の\R{符号}{ふ|ごう}だけがちがう\R{点}{てん}は、$y$ \R{軸}{じく}をはさんで\R{左右}{さ|ゆう}にちょうど\R{向}{む}かい\R{合}{あ}っているかな？
\end{enumerate}
\end{minipage}\hfill
\begin{minipage}[c]{0.36\linewidth}\centering
\begin{tikzpicture}[x=1mm, y=1mm]
  \bkthrcgrid{-4}{4}{-5}{6}{4.6}{0}
  \begin{scope}[x=4.6mm, y=4.6mm]
  \filldraw[fill=bkOrange!25, draw=bkOrange!80!black, line width=1.1pt, line join=round]
    (0,5) -- (2,2) -- (2,-2) -- (3,-4) -- (1,-3) -- (-1,-3) -- (-3,-4) -- (-2,-2) -- (-2,2) -- cycle;
  \foreach \p in {(0,5),(2,2),(2,-2),(3,-4),(1,-3),(-1,-3),(-3,-4),(-2,-2),(-2,2)} \fill[bkInk] \p circle (0.8mm);
  \end{scope}
\end{tikzpicture}
\end{minipage}
\par\vspace{3pt}
\R{左右}{さ|ゆう}がそろった\R{形}{かたち}になるのは、$(a,\ b)$ と $(-a,\ b)$ が、$y$ \R{軸}{じく}から\R{同}{おな}じ\R{距離}{きょ|り} $\lvert a\rvert$ だけ、\R{反対}{はん|たい}の\R{向}{む}きにはなれているからです。
\end{yattemiyou}

\begin{column}{\R{表計算}{ひょう|けい|さん}の「B3」も2つの\R{数}{かず}}
2-18 の\R{表計算}{ひょう|けい|さん}では、セルの\R{場所}{ば|しょ}を「B3」のような\R{番地}{ばん|ち}で\R{表}{あらわ}しました。B は\R{左}{ひだり}から2\R{番}{ばん}めの\R{列}{れつ}、3は\R{上}{うえ}から3\R{番}{ばん}めの\R{行}{ぎょう}です。これも、\R{場所}{ば|しょ}を「\R{横}{よこ}の\R{位置}{い|ち}」と「たての\R{位置}{い|ち}」の2つで\R{表}{あらわ}す\R{方法}{ほう|ほう}です。ただし\R{表計算}{ひょう|けい|さん}の\R{行}{ぎょう}は\R{下}{した}へ\R{行}{い}くほど\R{番号}{ばん|ごう}が\R{大}{おお}きくなり、\R{座標}{ざ|ひょう}の $y$（\R{上}{うえ}が＋）とは\R{向}{む}きが\R{反対}{はん|たい}です。\R{教室}{きょう|しつ}の\R{席}{せき}も「まどから4\R{列}{れつ}め、\R{前}{まえ}から3\R{番}{ばん}め」と、2つの\R{数}{かず}で1つに\R{決}{き}まります。\bkhatenaref の\R{答}{こた}えは「\R{数}{かず}は2つ」です。
\end{column}

\bksubsection{3}{マイナスの\R{座標}{ざ|ひょう}と、4つの\R{区画}{く|かく}}

\leavevmode\label{bk:3-2:kukaku}$x$ \R{軸}{じく}も $y$ \R{軸}{じく}も\R{数直線}{すう|ちょく|せん}なので、0の\R{左}{ひだり}や\R{下}{した}にはマイナスの\R{数}{かず}が\R{並}{なら}びます（2-2：\R{温度計}{おん|ど|けい}をたおすと\R{数直線}{すう|ちょく|せん}）。マイナスは「\R{反対}{はん|たい}の\R{向}{む}き」なので、$x$ \R{座標}{ざ|ひょう}が\R{負}{ふ}なら\R{左}{ひだり}へ、$y$ \R{座標}{ざ|ひょう}が\R{負}{ふ}なら\R{下}{した}へ\R{進}{すす}みます。

\noindent\begin{minipage}[c]{0.58\linewidth}\centering
\begin{tikzpicture}[x=1mm, y=1mm]
  \fill[bkOrange!8] (0,0) rectangle (32,24);
  \fill[bkBlue!8] (0,0) rectangle (-32,24);
  \fill[bkGreen!8] (0,0) rectangle (-32,-24);
  \fill[bkPurple!8] (0,0) rectangle (32,-24);
  \bkthrcgrid{-4}{4}{-3}{3}{8}{1}
  \draw[bkBlue, line width=1.1pt, -{Stealth[length=4pt]}] (-0.5,-1.2) -- (-23.5,-1.2);
  \draw[bkBlue, line width=1.1pt, -{Stealth[length=4pt]}] (-24,-0.8) -- (-24,14.5);
  \bkthrcpt{(-3,2)}{8}{B$(-3,\ 2)$}{south}{bkBlue}
  \bkthrcpt{(-2,-3)}{8}{C$(-2,\ -3)$}{north}{bkGreen!80!black}
  \bkthrcpt{(4,-1)}{8}{D$(4,\ -1)$}{north east}{bkPurple}
  \bkthrcpt{(3,2)}{8}{A$(3,\ 2)$}{south}{bkRed}
\end{tikzpicture}
\end{minipage}\hfill
\begin{minipage}[c]{0.38\linewidth}
\R{点}{てん} B$(-3,\ 2)$ は、O から\R{左}{ひだり}へ3、そこから\R{上}{うえ}へ2（\R{青}{あお}い\R{矢印}{や|じるし}）。\par\vspace{2pt}
C$(-2,\ -3)$ は、\R{左}{ひだり}へ2、\R{下}{した}へ3。\par\vspace{2pt}
D$(4,\ -1)$ は、\R{右}{みぎ}へ4、\R{下}{した}へ1。
\end{minipage}

\par\vspace{2mm}
2\R{本}{ほん}の\R{軸}{じく}で、\R{平面}{へい|めん}は4つの\R{区画}{く|かく}に\R{分}{わ}かれます。どの\R{区画}{く|かく}にあるかは、$x$ \R{座標}{ざ|ひょう}と $y$ \R{座標}{ざ|ひょう}の\R{符号}{ふ|ごう}（＋か－か）だけで\R{決}{き}まります。

\begin{bkcenter}
\begin{bktable}{colspec={Q[l,m,wd=24mm] Q[c,m,wd=18mm] Q[c,m,wd=18mm] Q[c,m,wd=30mm]}, rowsep=0.8pt, cells={font=\small}, row{1}={font=\sffamily\bfseries\small}}
\R{区画}{く|かく} & $x$ \R{座標}{ざ|ひょう} & $y$ \R{座標}{ざ|ひょう} & \R{例}{れい} \\
\SetCell{bg=bkOrange!8} \R{右上}{みぎ|うえ} & ＋ & ＋ & A$(3,\ 2)$ \\
\SetCell{bg=bkBlue!8} \R{左上}{ひだり|うえ} & － & ＋ & B$(-3,\ 2)$ \\
\SetCell{bg=bkGreen!8} \R{左下}{ひだり|した} & － & － & C$(-2,\ -3)$ \\
\SetCell{bg=bkPurple!8} \R{右下}{みぎ|した} & ＋ & － & D$(4,\ -1)$ \\
\end{bktable}
\bktabcaption{\R{符号}{ふ|ごう}で、\R{区画}{く|かく}が\R{決}{き}まる。}
\end{bkcenter}

$(4,\ 0)$ や $(0,\ -3)$ のように、\R{座標}{ざ|ひょう}のどちらかが0の\R{点}{てん}は\R{軸}{じく}の\R{上}{うえ}にあり、どの\R{区画}{く|かく}にも\R{入}{はい}りません（\R{高校}{こう|こう}では、4つの\R{区画}{く|かく}を\R{象限}{しょう|げん}とよび、\R{右上}{みぎ|うえ}から\R{反時計回}{はん|とけい|まわ}りに\R{第}{だい}1\R{象限}{しょう|げん}、\R{第}{だい}2\R{象限}{しょう|げん}、…と\R{名前}{な|まえ}をつけます）。

\bksubsection{4}{＋の\R{向}{む}きと\R{原点}{げん|てん}は、\R{自分}{じ|ぶん}で\R{決}{き}める}

\leavevmode\label{bk:3-2:kimeru}\R{座標}{ざ|ひょう}を\R{使}{つか}うとき、\R{原点}{げん|てん}をどこに\R{置}{お}くか、どちらを＋の\R{向}{む}きにするかは、\R{問題}{もん|だい}を\R{解}{と}く\R{人}{ひと}が\R{自分}{じ|ぶん}で\R{決}{き}めてかまいません。\R{宝}{たから}さがしでは\R{木}{き}を\R{原点}{げん|てん}にしましたが、\R{池}{いけ}を\R{原点}{げん|てん}にしてもよかったのです。ただし、\R{大事}{だい|じ}な\R{約束}{やく|そく}が2つあります。

\begin{taisetsu}
① \R{原点}{げん|てん}と＋の\R{向}{む}きを\R{決}{き}めたら、\R{図}{ず}にかき、\R{問題}{もん|だい}の\R{途中}{と|ちゅう}で\R{変}{か}えない。\par
② \R{原点}{げん|てん}は、\R{計算}{けい|さん}が\R{楽}{らく}になる\R{所}{ところ}に\R{置}{お}く。
\end{taisetsu}

\noindent\begin{minipage}[c]{0.55\linewidth}
\R{手}{て}で\R{真上}{ま|うえ}にボールを\R{投}{な}げます。ボールをはなした\R{手}{て}は\R{地面}{じ|めん}から\qty{1.5}{m}、ボールのいちばん\R{高}{たか}い\R{所}{ところ}は\R{地面}{じ|めん}から\qty{6.5}{m}でした。\R{上}{うえ}を＋にした $y$ \R{軸}{じく}で、\R{原点}{げん|てん}の\R{置}{お}き\R{方}{かた}を2\R{通}{とお}り\R{考}{かんが}えます（\R{右}{みぎ}の\R{図}{ず}の\R{赤}{あか}い\R{軸}{じく}と\R{青}{あお}い\R{軸}{じく}）。
\end{minipage}\hfill
\begin{minipage}[c]{0.41\linewidth}\centering
\begin{tikzpicture}[x=1mm, y=1mm]
  \fill[bkGreen!25] (-4,-2) rectangle (50,0);
  \draw[bkGreen!60!black, line width=0.8pt] (-4,0) -- (50,0);
  \bkkid{(14,0)}{16}{bkBlue}{bkHair}
  \draw[bkGray, line width=0.6pt, dash pattern=on 2pt off 1.5pt] (20,13.5) -- (20,46.2);
  \shade[ball color=bkOrange] (20,48) circle (1.8);
  \shade[ball color=bkOrange] (20,12) circle (1.5);
  % 左の軸：地面が原点
  \draw[-{Stealth[length=3pt]}, bkRed, line width=0.8pt] (4,0) -- (4,54);
  \foreach \v/\yy in {0/0, 1.5/12, 6.5/48} {\draw[bkRed, line width=0.6pt] (3,\yy) -- (5,\yy);
    \node[anchor=east, font=\footnotesize, text=bkRed!80!black, inner sep=1pt] at (2.5,\yy) {$\v$};}
  % 右の軸：手が原点
  \draw[-{Stealth[length=3pt]}, bkBlue, line width=0.8pt] (36,-1) -- (36,54);
  \foreach \v/\yy in {-1.5/0, 0/12, 5.0/48} {\draw[bkBlue, line width=0.6pt] (35,\yy) -- (37,\yy);
    \node[anchor=west, font=\footnotesize, text=bkBlue!80!black, inner sep=1pt] at (37.5,\yy+1.6) {$\v$};}
  \draw[bkGray, line width=0.4pt, densely dotted] (5,12) -- (35,12) (5,48) -- (35,48);
\end{tikzpicture}
\end{minipage}

\par\vspace{2mm}
\textbf{\R{地面}{じ|めん}を\R{原点}{げん|てん}にする（\R{赤}{あか}）}：\R{手}{て}は $y=1.5$、いちばん\R{高}{たか}い\R{所}{ところ}は $y=6.5$。\R{上}{あ}がった\R{高}{たか}さは $6.5-1.5=5.0$（m）。\par\vspace{2pt}
\textbf{\R{手}{て}を\R{原点}{げん|てん}にする（\R{青}{あお}）}：\R{手}{て}は $y=0$、いちばん\R{高}{たか}い\R{所}{ところ}は $y=5.0$。\R{上}{あ}がった\R{高}{たか}さは $5.0-0=5.0$（m）。\par\vspace{2pt}
\R{原点}{げん|てん}を\R{変}{か}えると、\R{座標}{ざ|ひょう}の\R{数}{かず}は\R{変}{か}わります。でも、2つの\R{点}{てん}のあいだの\R{長}{なが}さ（\R{上}{あ}がった\R{高}{たか}さ）は、どちらでも\qty{5.0}{m}で\R{同}{おな}じです。\R{原点}{げん|てん}は「ものさしの0をどこに\R{当}{あ}てるか」を\R{決}{き}めているだけだからです。「どれだけ\R{上}{あ}がったか」を\R{知}{し}りたいなら、はなした\R{所}{ところ}を\R{原点}{げん|てん}にすると、ひき\R{算}{ざん}がいらず\R{楽}{らく}です。

\bkthrchead{\R{途中}{と|ちゅう}で\R{原点}{げん|てん}を\R{変}{か}えると}
\R{手}{て}の\R{高}{たか}さは\R{地面}{じ|めん}を\R{原点}{げん|てん}にして1.5、いちばん\R{高}{たか}い\R{所}{ところ}は\R{手}{て}を\R{原点}{げん|てん}にして5.0と、ばらばらに\R{決}{き}めてしまうと、$5.0-1.5=3.5$ となって、まちがった\R{答}{こた}えが\R{出}{で}ます。ほんとうは\qty{5.0}{m}\R{上}{あ}がっています。①の\R{約束}{やく|そく}は、このまちがいをふせぐためです。

\bkthrchead{\R{下}{した}を＋にしてもよい}
\R{落}{お}ちる\R{物}{もの}だけを\R{考}{かんが}えるときは、\R{下}{した}を＋にすると、\R{座標}{ざ|ひょう}がマイナスにならず\R{計算}{けい|さん}が\R{楽}{らく}です。たとえば、はなしてから1\R{秒}{びょう}で\qty{4.9}{m}\R{落}{お}ちたボールは、\R{下}{した}を＋にすれば $y=4.9$、\R{上}{うえ}を＋にすれば $y=-4.9$ です。どちらも\R{正}{ただ}しい\R{表}{あらわ}し\R{方}{かた}です。コンピュータの\R{画面}{が|めん}では、\R{原点}{げん|てん}を\R{左上}{ひだり|うえ}のすみに\R{置}{お}き、$y$ を\R{下}{した}\R{向}{む}きに＋にするのがふつうです（\pageref{bk:3-2:gamen}ページ）。

\bksubsection{5}{\R{点}{てん}を\R{動}{うご}かす}

\noindent\label{bk:3-2:ugokasu}\begin{minipage}[c]{0.44\linewidth}
ゲームで、キャラクターが\R{方眼}{ほう|がん}の\R{上}{うえ}を\R{動}{うご}きます。いまの\R{場所}{ば|しょ}は A$(2,\ 1)$ です。\R{右}{みぎ}へ3マス\R{動}{うご}くと、どこへ\R{行}{い}くでしょう。

\R{右}{みぎ}へ\R{動}{うご}くので、$x$ \R{座標}{ざ|ひょう}が3ふえて $2+3=5$。\R{上下}{じょう|げ}には\R{動}{うご}かないので、$y$ \R{座標}{ざ|ひょう}は1のまま。B$(5,\ 1)$ です。そこから\R{上}{うえ}へ2マス\R{動}{うご}くと、$y$ \R{座標}{ざ|ひょう}が2ふえて $1+2=3$。C$(5,\ 3)$ です。
\end{minipage}\hfill
\begin{minipage}[c]{0.52\linewidth}\centering
\begin{tikzpicture}[x=1mm, y=1mm]
  \bkthrcgrid{-1}{6}{-1}{4}{8}{1}
  \draw[bkOrange, line width=1.3pt, -{Stealth[length=5pt]}] (17.5,8) -- (38.5,8);
  \draw[bkOrange, line width=1.3pt, -{Stealth[length=5pt]}] (40,9.5) -- (40,22.5);
  \node[anchor=north, font=\sffamily\footnotesize, text=bkOrange!80!black, fill=white, inner sep=0.6pt] at (28,6.6) {\R{右}{みぎ}へ3};
  \node[anchor=west, font=\sffamily\footnotesize, text=bkOrange!80!black, fill=white, inner sep=0.6pt] at (41.5,16) {\R{上}{うえ}へ2};
  \bkthrcpt{(2,1)}{8}{A}{south}{bkBlue}
  \bkthrcpt{(5,1)}{8}{B}{north west}{bkBlue}
  \bkthrcpt{(5,3)}{8}{C}{south}{bkRed}
\end{tikzpicture}
\end{minipage}

\par\vspace{2mm}
\R{数}{かず}だけで\R{書}{か}くと、\R{次}{つぎ}のようになります。

\begin{bkcenter}
\begin{bktable}{colspec={Q[l,m,wd=34mm] Q[c,m,wd=26mm] Q[c,m,wd=26mm]}, rowsep=0.8pt, cells={font=\small}, row{1}={font=\sffamily\bfseries\small}}
\R{動}{うご}き & $x$ \R{座標}{ざ|ひょう} & $y$ \R{座標}{ざ|ひょう} \\
\R{右}{みぎ}へ $p$ & $p$ ふえる（$+p$） & そのまま \\
\R{左}{ひだり}へ $p$ & $p$ へる（$-p$） & そのまま \\
\R{上}{うえ}へ $q$ & そのまま & $q$ ふえる（$+q$） \\
\R{下}{した}へ $q$ & そのまま & $q$ へる（$-q$） \\
\end{bktable}
\bktabcaption{\R{点}{てん} $(a,\ b)$ を\R{動}{うご}かすと}
\end{bkcenter}

\R{点}{てん} $(a,\ b)$ を\R{右}{みぎ}へ3\R{動}{うご}かすと $(a+3,\ b)$、\R{左}{ひだり}へ2\R{動}{うご}かすと $(a-2,\ b)$、\R{右}{みぎ}へ $p$・\R{上}{うえ}へ $q$ \R{動}{うご}かすと $(a+p,\ b+q)$ です。\R{横}{よこ}の\R{動}{うご}きは $x$ \R{座標}{ざ|ひょう}だけ、たての\R{動}{うご}きは $y$ \R{座標}{ざ|ひょう}だけを\R{変}{か}えます。\R{横}{よこ}とたては、おたがいにじゃまをしません。

\bkthrchead{どれだけ\R{動}{うご}いたかを\R{求}{もと}める}
\R{反対}{はん|たい}に、A$(2,\ 1)$ から C$(5,\ 3)$ へ\R{動}{うご}いたとき、どれだけ\R{動}{うご}いたかは、「あと－まえ」（2-3）で\R{求}{もと}めます。$x$ は $5-2=3$ で\R{右}{みぎ}へ3、$y$ は $3-1=2$ で\R{上}{うえ}へ2です。ひき\R{算}{ざん}の\R{答}{こた}えがマイナスなら、\R{左}{ひだり}や\R{下}{した}へ\R{動}{うご}いたということです。

\bksubsection{6}{\R{点}{てん}がたくさんあるとき：\R{下}{した}つき\R{文字}{も|じ}}

\noindent\label{bk:3-2:shitatsuki}\begin{minipage}[c]{0.48\linewidth}
\R{部屋}{へ|や}のそうじをするロボットそうじ\R{機}{き}は、\R{自分}{じ|ぶん}の\R{場所}{ば|しょ}を\R{座標}{ざ|ひょう}で\R{覚}{おぼ}えながら\R{動}{うご}きます。1\R{秒}{びょう}ごとの\R{場所}{ば|しょ}を\R{記録}{き|ろく}すると、\R{座標}{ざ|ひょう}がたくさん\R{並}{なら}びます。

1\R{回}{かい}めに\R{記録}{き|ろく}した\R{座標}{ざ|ひょう}を $(x_1,\ y_1)$、2\R{回}{かい}めを $(x_2,\ y_2)$、3\R{回}{かい}めを $(x_3,\ y_3)$ と\R{書}{か}きます。$x$ の\R{右下}{みぎ|した}の\R{小}{ちい}さな\R{数字}{すう|じ}は「\R{何番}{なん|ばん}めか」を\R{表}{あらわ}す\R{名札}{な|ふだ}で、これを\textbf{\R{下}{した}つき\R{文字}{も|じ}}といいます。\bkidx{したつきもじ}{\R{下}{した}つき\R{文字}{も|じ}}$x_1$ は「エックス・いち」、$x_2$ は「エックス・に」と\R{読}{よ}みます。
\end{minipage}\hfill
\begin{minipage}[c]{0.48\linewidth}\centering
\begin{tikzpicture}[x=1mm, y=1mm]
  \bkthrcgrid{-1}{6}{-1}{4}{7.5}{1}
  \draw[bkGray, line width=0.8pt, densely dashed] (7.5,7.5) -- (22.5,22.5) -- (37.5,15);
  \bkthrcpt{(1,1)}{7.5}{$(x_1,\ y_1)$}{north west}{bkBlue}
  \bkthrcpt{(3,3)}{7.5}{$(x_2,\ y_2)$}{south}{bkBlue}
  \bkthrcpt{(5,2)}{7.5}{$(x_3,\ y_3)$}{north}{bkBlue}
\end{tikzpicture}
\end{minipage}

\par\vspace{2mm}
\R{図}{ず}の\R{例}{れい}では、$(x_1,\ y_1)=(1,\ 1)$、$(x_2,\ y_2)=(3,\ 3)$、$(x_3,\ y_3)=(5,\ 2)$ です。$x_2=3$、$y_3=2$ のように、1つずつ\R{取}{と}り\R{出}{だ}して\R{使}{つか}うこともできます。\R{点}{てん}が100\R{個}{こ}あっても、$a$、$b$、$c$、…のように\R{文字}{も|じ}が\R{足}{た}りなくなることはありません。

\begin{bkcenter}
\begin{bktable}{colspec={Q[c,m,wd=18mm] Q[l,m,wd=34mm] Q[l,m,wd=56mm]}, rowsep=0.8pt, cells={font=\small}, row{1}={font=\sffamily\bfseries\small}}
\R{書}{か}き\R{方}{かた} & \R{読}{よ}み\R{方}{かた} & \R{意味}{い|み} \\
$x_1$ & エックス・いち & 1\R{番}{ばん}めの $x$（\R{名札}{な|ふだ}） \\
$x^2$ & エックスの2\R{乗}{じょう} & $x\times x$（2-7） \\
$2x$ & 2エックス & $2\times x$（2-12） \\
\end{bktable}
\bktabcaption{\R{小}{ちい}さな\R{数字}{すう|じ}の\R{位置}{い|ち}で、\R{意味}{い|み}がまったくちがう。}
\end{bkcenter}

\R{下}{した}つき\R{文字}{も|じ}は、かけ\R{算}{ざん}ではありません。$x_1$ を $x\times1$ と\R{読}{よ}むと、ただの $x$ になってしまい、\R{名札}{な|ふだ}の\R{意味}{い|み}がなくなります。2\R{回}{かい}めと1\R{回}{かい}めのあいだに $x$ が\R{変}{か}わった\R{量}{りょう}は、「あと－まえ」で $x_2-x_1$ です。2-3 の $\Delta x$ と\R{同}{おな}じものです。\R{図}{ず}の\R{例}{れい}では $x_2-x_1=3-1=2$ で、ロボットは\R{右}{みぎ}へ2\R{動}{うご}きました。3-13 からのデータの\R{章}{しょう}でも、1\R{回}{かい}めにはかった\R{値}{あたい}を $x_1$、2\R{回}{かい}めを $x_2$ と\R{書}{か}いて\R{使}{つか}います。

\bksubsection{7}{0からの\R{距離}{きょ|り}：\R{絶対値}{ぜっ|たい|ち}と、2\R{点}{てん}の\R{距離}{きょ|り}}

\label{bk:3-2:kyori}\bkthrchead{\R{数直線}{すう|ちょく|せん}で\R{見}{み}る}
\noindent\begin{minipage}[c]{0.46\linewidth}
\R{数直線}{すう|ちょく|せん}の $-3$ と3は、0の\R{左}{ひだり}と\R{右}{みぎ}にあり、\R{向}{む}きは\R{反対}{はん|たい}です。でも、0からどれだけはなれているかは、どちらも3\R{目}{め}もりで\R{同}{おな}じです。\R{海面}{かい|めん}から\R{上}{うえ}へ\qty{50}{m}の\R{山}{やま}の\R{上}{うえ}と、\R{下}{した}へ\qty{50}{m}の\R{海}{うみ}の\R{中}{なか}は、\R{向}{む}きはちがっても、\R{海面}{かい|めん}からの\R{距離}{きょ|り}は\R{同}{おな}じ\qty{50}{m}です。
\end{minipage}\hfill
\begin{minipage}[c]{0.50\linewidth}\centering
\begin{tikzpicture}[x=1mm, y=1mm]
  \draw[bkInk, line width=0.8pt, -{Stealth[length=4pt]}] (-31,0) -- (32,0);
  \foreach \k in {-4,...,4} {\draw[bkInk, line width=0.5pt] (\k*7,-1.2) -- (\k*7,1.2);
    \node[anchor=north, font=\footnotesize, inner sep=0pt] at (\k*7,-2.4) {$\k$};}
  \draw[bkOrange, line width=1.6pt] (-21,3.5) -- (0,3.5);
  \draw[bkBlue, line width=1.6pt] (0,3.5) -- (21,3.5);
  \fill[bkOrange] (-21,0) circle (1.3);
  \fill[bkBlue] (21,0) circle (1.3);
  \node[anchor=south, font=\sffamily\footnotesize, text=bkOrange!80!black] at (-10.5,4.5) {3\R{目}{め}もり};
  \node[anchor=south, font=\sffamily\footnotesize, text=bkBlue!80!black] at (10.5,4.5) {3\R{目}{め}もり};
\end{tikzpicture}
\end{minipage}

\par\vspace{2mm}
\R{数直線}{すう|ちょく|せん}の\R{上}{うえ}で、ある\R{数}{かず}が0からどれだけはなれているかを、その\R{数}{かず}の\textbf{\R{絶対値}{ぜっ|たい|ち}}といいます。\bkidx{ぜったいち}{\R{絶対値}{ぜっ|たい|ち}}$-3$ の\R{絶対値}{ぜっ|たい|ち}を、\R{数}{かず}を2\R{本}{ほん}のたての\R{線}{せん}ではさんで $\lvert -3\rvert $ と\R{書}{か}き、「\R{絶対値}{ぜっ|たい|ち}マイナス3」と\R{読}{よ}みます。
\[ \lvert -3\rvert =3,\qquad \lvert 3\rvert =3,\qquad \lvert 0\rvert =0,\qquad \lvert -2.5\rvert =2.5 \]

\bkthrchead{\R{記号}{き|ごう}で\R{見}{み}る}
\R{上}{うえ}の\R{例}{れい}を\R{見}{み}ると、\R{絶対値}{ぜっ|たい|ち}は「マイナスの\R{符号}{ふ|ごう}をとった\R{数}{かず}」になっています。\R{正}{せい}の\R{数}{かず}と0は、そのままです。\R{距離}{きょ|り}はマイナスにならないので、\R{絶対値}{ぜっ|たい|ち}はいつも0か\R{正}{せい}の\R{数}{かず}です。\R{距離}{きょ|り}で\R{考}{かんが}えても、\R{符号}{ふ|ごう}をとると\R{考}{かんが}えても、\R{同}{おな}じ\R{答}{こた}えになります。

\bkthrchead{\R{数直線}{すう|ちょく|せん}の2\R{点}{てん}の\R{距離}{きょ|り}}
\noindent\begin{minipage}[c]{0.44\linewidth}
\R{朝}{あさ}の\R{気温}{き|おん}が\qty{-2}{\degreeCelsius}、\R{昼}{ひる}が\qty{3}{\degreeCelsius}でした。\R{温度計}{おん|ど|けい}の\R{目}{め}もりを\R{数}{かぞ}えると、$-2$ から3までは5\R{目}{め}もりです。\R{計算}{けい|さん}では「\R{大}{おお}きいほう－\R{小}{ちい}さいほう」で $3-(-2)=3+2=5$ です（2-2：ひく\R{数}{かず}がマイナスなら、たす）。

\end{minipage}\hfill
\begin{minipage}[c]{0.52\linewidth}\centering
\begin{tikzpicture}[x=1mm, y=1mm]
  \draw[bkInk, line width=0.8pt, -{Stealth[length=4pt]}] (-27,0) -- (36,0);
  \foreach \k in {-3,...,4} {\draw[bkInk, line width=0.5pt] (\k*8,-1.2) -- (\k*8,1.2);
    \node[anchor=north, font=\footnotesize, inner sep=0pt] at (\k*8,-2.4) {$\k$};}
  \fill[bkBlue] (-16,0) circle (1.3);
  \fill[bkRed] (24,0) circle (1.3);
  \draw[{Stealth[length=3pt]}-{Stealth[length=3pt]}, bkGreen!70!black, line width=0.9pt] (-16,4) -- (24,4);
  \node[anchor=south, font=\sffamily\footnotesize, text=bkGreen!50!black] at (4,5) {5\R{目}{め}もり};
  \node[anchor=north, font=\sffamily\footnotesize, text=bkBlue!80!black] at (-16,-6.5) {\R{朝}{あさ}};
  \node[anchor=north, font=\sffamily\footnotesize, text=bkRed!80!black] at (24,-6.5) {\R{昼}{ひる}};
\end{tikzpicture}
\end{minipage}

\par\vspace{2mm}
\R{順番}{じゅん|ばん}を\R{逆}{ぎゃく}にひくと $-2-3=-5$ になりますが、\R{絶対値}{ぜっ|たい|ち}をとれば $\lvert -5\rvert =5$ で、\R{同}{おな}じ\R{距離}{きょ|り}になります。だから、どちらからひいてもよいように、2\R{点}{てん} $a$、$b$ の\R{距離}{きょ|り}は $\lvert b-a\rvert $ と\R{書}{か}きます。

\bkthrchead{\R{横}{よこ}・たてに\R{並}{なら}んだ2\R{点}{てん}の\R{距離}{きょ|り}}
\noindent\begin{minipage}[c]{0.56\linewidth}\centering
\begin{tikzpicture}[x=1mm, y=1mm]
  \bkthrcgrid{-3}{4}{-3}{3}{8}{1}
  \draw[{Stealth[length=3pt]}-{Stealth[length=3pt]}, bkGreen!70!black, line width=1.1pt] (-14.5,8) -- (22.5,8);
  \draw[{Stealth[length=3pt]}-{Stealth[length=3pt]}, bkPurple, line width=1.1pt] (24,6.5) -- (24,-14.5);
  \draw[bkGray, line width=0.8pt, dash pattern=on 2pt off 1.5pt] (-16,8) -- (24,-16);
  \node[font=\sffamily\bfseries\small, text=bkGray, fill=white, inner sep=0.8pt] at (16,-6.5) {?};
  \bkthrcpt{(-2,1)}{8}{A}{south}{bkBlue}
  \bkthrcpt{(3,1)}{8}{B}{south}{bkBlue}
  \bkthrcpt{(3,-2)}{8}{C}{north}{bkBlue}
\end{tikzpicture}
\end{minipage}\hfill
\begin{minipage}[c]{0.40\linewidth}
A$(-2,\ 1)$ と B$(3,\ 1)$ は、$y$ \R{座標}{ざ|ひょう}がどちらも1なので、\R{横}{よこ}に\R{並}{なら}んでいます。\R{距離}{きょ|り}は $x$ \R{座標}{ざ|ひょう}だけを\R{比}{くら}べて、$\lvert 3-(-2)\rvert =5$ です。マス\R{目}{め}を\R{数}{かぞ}えても5です。\par\vspace{2pt}
B$(3,\ 1)$ と C$(3,\ -2)$ は、$x$ \R{座標}{ざ|ひょう}が\R{同}{おな}じで、たてに\R{並}{なら}んでいます。\R{距離}{きょ|り}は $\lvert -2-1\rvert =3$ です。
\end{minipage}

\par\vspace{2mm}
\R{下}{した}つき\R{文字}{も|じ}で\R{書}{か}くと、\R{横}{よこ}に\R{並}{なら}んだ2\R{点}{てん} $(x_1,\ y_1)$、$(x_2,\ y_1)$ の\R{距離}{きょ|り}は $\lvert x_2-x_1\rvert $、たてに\R{並}{なら}んだ2\R{点}{てん} $(x_1,\ y_1)$、$(x_1,\ y_2)$ の\R{距離}{きょ|り}は $\lvert y_2-y_1\rvert $ です。

A と C のように\R{斜}{なな}めに\R{並}{なら}んだ2\R{点}{てん}の\R{距離}{きょ|り}は、5と3をたした8ではありません（\R{斜}{なな}めの\R{近道}{ちか|みち}のほうが\R{短}{みじか}いからです）。\R{斜}{なな}めの\R{距離}{きょ|り}は、\R{第}{だい}4\R{巻}{かん}の 4-4（\R{三平方}{さん|へい|ほう}の\R{定理}{てい|り}）で\R{求}{もと}められるようになります。

\begin{motto}[2]
\textbf{$y=\lvert x\rvert $ の\R{点}{てん}をうってみよう}\par
$x$ が $-3$、$-2$、…、3のときの $y=\lvert x\rvert $ を\R{表}{ひょう}にして、\R{点}{てん}をうつと、\R{下}{した}のようにV\R{字}{じ}の\R{形}{かたち}に\R{並}{なら}びます。\R{原点}{げん|てん}で、\R{線}{せん}がとがって\R{折}{お}れ\R{曲}{ま}がっています。このとがった\R{所}{ところ}は、\R{第}{だい}5\R{巻}{かん}の 5-3 で、もう\R{一度}{いち|ど}\R{調}{しら}べます。\par\vspace{2pt}
\noindent\begin{minipage}[c]{0.52\linewidth}
\begin{bktable}{colspec={Q[c,m,wd=7mm] *{7}{Q[c,m,wd=5.6mm]}}, rowsep=0.5pt, colsep=1.6pt, cells={font=\small}}
$x$ & $-3$ & $-2$ & $-1$ & 0 & 1 & 2 & 3 \\
$\lvert x\rvert $ & 3 & 2 & 1 & 0 & 1 & 2 & 3 \\
\end{bktable}
\end{minipage}\hfill
\begin{minipage}[c]{0.44\linewidth}\centering
\begin{tikzpicture}[x=1mm, y=1mm]
  \bkthrcgrid{-3}{3}{0}{3}{6}{1}
  \draw[bkBlue, line width=1pt] (-18,18) -- (0,0) -- (18,18);
  \foreach \k in {-3,...,3} {\pgfmathtruncatemacro\a{abs(\k)}\fill[bkInk] (\k*6,\a*6) circle (1);}
\end{tikzpicture}
\end{minipage}
\end{motto}

\begin{kotoba}
\small
\bkword{\R{座標軸}{ざ|ひょう|じく}}{\R{横}{よこ}の $x$ \R{軸}{じく}とたての $y$ \R{軸}{じく}。0の\R{所}{ところ}で\R{直角}{ちょっ|かく}に\R{交}{まじ}わる。}
\bkword{\R{原点}{げん|てん}}{2\R{本}{ほん}の\R{軸}{じく}が\R{交}{まじ}わる\R{点}{てん}。O と\R{書}{か}く。\R{座標}{ざ|ひょう}は $(0,\ 0)$。}
\bkword{\R{座標}{ざ|ひょう}}{\R{点}{てん}の\R{場所}{ば|しょ}を\R{表}{あらわ}す\R{数}{かず}の\R{組}{くみ}。$(x\text{\R{座標}{ざ|ひょう}},\ y\text{\R{座標}{ざ|ひょう}})$ の\R{順}{じゅん}に\R{書}{か}く。}
\bkword{\R{下}{した}つき\R{文字}{も|じ}}{$x_1$、$x_2$ のように、\R{何番}{なん|ばん}めかを\R{表}{あらわ}す\R{右下}{みぎ|した}の\R{小}{ちい}さな\R{数字}{すう|じ}。かけ\R{算}{ざん}ではない。}
\bkword{\R{絶対値}{ぜっ|たい|ち}}{\R{数直線}{すう|ちょく|せん}で0からの\R{距離}{きょ|り}。$\lvert -3\rvert =3$。}
\end{kotoba}

\bkkigocards{\bkkigocard{$x_1$}{エックス・いち}{1\R{番}{ばん}めの $x$。\R{右下}{みぎ|した}の\R{数字}{すう|じ}は\R{名札}{な|ふだ}で、かけ\R{算}{ざん}ではない}{x sub one}{3-2}}{\bkkigocard{$\lvert x\rvert $}{\R{絶対値}{ぜっ|たい|ち}エックス}{$x$ の\R{絶対値}{ぜっ|たい|ち}。\R{数直線}{すう|ちょく|せん}で0からの\R{距離}{きょ|り}。$\lvert -3\rvert =3$}{absolute value}{3-2}}

\begin{column}{\R{将棋}{しょう|ぎ}の「7\R{六}{ろく}\R{歩}{ふ}」も、\R{地球}{ち|きゅう}の\R{緯度}{い|ど}・\R{経度}{けい|ど}も\R{座標}{ざ|ひょう}}
\R{将棋}{しょう|ぎ}の\R{記録}{き|ろく}では、「7\R{六}{ろく}\R{歩}{ふ}」のように、\R{盤}{ばん}の\R{右}{みぎ}から\R{数}{かぞ}えた\R{列}{れつ}（1〜9の\R{数字}{すう|じ}）と、\R{上}{うえ}から\R{数}{かぞ}えた\R{段}{だん}（\R{一}{いち}〜\R{九}{きゅう}の\R{漢数字}{かん|すう|じ}）の2つで\R{場所}{ば|しょ}を\R{表}{あらわ}します。\R{地球}{ち|きゅう}の\R{上}{うえ}の\R{場所}{ば|しょ}は、\R{赤道}{せき|どう}から\R{北}{きた}や\R{南}{みなみ}へ\R{何度}{なん|ど}かを\R{表}{あらわ}す\R{緯度}{い|ど}と、イギリスのグリニッジから\R{東}{ひがし}や\R{西}{にし}へ\R{何度}{なん|ど}かを\R{表}{あらわ}す\R{経度}{けい|ど}の2つで\R{表}{あらわ}します。たとえば\R{東京駅}{とう|きょう|えき}は、\R{北緯}{ほく|い}\R{約}{やく}35.7\R{度}{ど}、\R{東経}{とう|けい}\R{約}{やく}139.8\R{度}{ど}です。スマホの\R{地図}{ち|ず}は、この2つの\R{数}{かず}で\R{今}{いま}いる\R{場所}{ば|しょ}を\R{示}{しめ}しています。どれも、「どこから・どちら\R{向}{む}きに・いくつ」を2\R{組}{くみ}\R{決}{き}めて、\R{場所}{ば|しょ}を2つの\R{数}{かず}で\R{伝}{つた}える\R{方法}{ほう|ほう}です。
\end{column}

\begin{taisetsu}
\R{点}{てん} $(a,\ b)$：\R{原点}{げん|てん}から\R{右}{みぎ}へ $a$、\R{上}{うえ}へ $b$（マイナスなら\R{左}{ひだり}・\R{下}{した}）。$x$ が\R{先}{さき}、$y$ が\R{後}{あと}。\par\vspace{2pt}
\R{右}{みぎ}へ $p$、\R{上}{うえ}へ $q$ \R{動}{うご}かすと $(a+p,\ b+q)$。\R{動}{うご}いた\R{量}{りょう}は $x_2-x_1$、$y_2-y_1$（あと－まえ）。\par\vspace{2pt}
\R{横}{よこ}に\R{並}{なら}んだ2\R{点}{てん}の\R{距離}{きょ|り}は $\lvert x_2-x_1\rvert$、たてなら $\lvert y_2-y_1\rvert$。
\end{taisetsu}

%% ---------------- 例題 ----------------
\bksection{reidai}{\R{例題}{れい|だい}でたしかめよう}

\begin{reidai}[\R{点}{てん}の\R{座標}{ざ|ひょう}を\R{読}{よ}む・\R{点}{てん}をかく]
\noindent\begin{minipage}[c]{0.46\linewidth}
（1）\R{右}{みぎ}の\R{図}{ず}の\R{点}{てん} A、B、C、D の\R{座標}{ざ|ひょう}を\R{答}{こた}えましょう。\par
（2）\R{点}{てん} E$(-4,\ 1)$ をかくには、\R{原点}{げん|てん}からどう\R{進}{すす}めばよいですか。\par
\kangaekata \R{点}{てん}から $x$ \R{軸}{じく}・$y$ \R{軸}{じく}へまっすぐな\R{線}{せん}をおろして、\R{当}{あ}たった\R{目}{め}もりを\R{読}{よ}みます。$x$ が\R{先}{さき}、$y$ が\R{後}{あと}です。
\end{minipage}\hfill
\begin{minipage}[c]{0.50\linewidth}\centering
\begin{tikzpicture}[x=1mm, y=1mm]
  \bkthrcgrid{-5}{5}{-4}{4}{4.6}{2}
  \bkthrcpt{(4,3)}{4.6}{A}{west}{bkRed}
  \bkthrcpt{(-2,4)}{4.6}{B}{west}{bkBlue}
  \bkthrcpt{(-3,-2)}{4.6}{C}{east}{bkGreen!80!black}
  \bkthrcpt{(2,-4)}{4.6}{D}{west}{bkPurple}
\end{tikzpicture}
\end{minipage}
\par
\shiki（1）\nobreak A：\R{下}{した}へおろすと $x$ \R{軸}{じく}の4、\R{横}{よこ}へ\R{引}{ひ}くと $y$ \R{軸}{じく}の3\par
\hspace*{2.6\zw}B：$x$ \R{軸}{じく}の $-2$、$y$ \R{軸}{じく}の4。C：$-3$ と $-2$。D：2と $-4$\par
\hspace*{2.6\zw}（2）\nobreak $x=-4$ はマイナスなので\R{左}{ひだり}へ4、$y=1$ はプラスなので\R{上}{うえ}へ1\par
\kotae（1）\nobreak A$(4,\ 3)$、B$(-2,\ 4)$、C$(-3,\ -2)$、D$(2,\ -4)$\par
\hspace*{2.6\zw}（2）\nobreak \R{原点}{げん|てん}から\R{左}{ひだり}へ4、そこから\R{上}{うえ}へ1\R{進}{すす}む
\end{reidai}

\begin{jibun}
\noindent\begin{minipage}[c]{0.46\linewidth}
（1）\R{右}{みぎ}の\R{図}{ず}の\R{点}{てん} P、Q、R、S の\R{座標}{ざ|ひょう}を\R{答}{こた}えましょう。\par
（2）\R{点}{てん} T$(3,\ -2)$ をかくには、\R{原点}{げん|てん}からどう\R{進}{すす}めばよいですか。
\end{minipage}\hfill
\begin{minipage}[c]{0.50\linewidth}\centering
\begin{tikzpicture}[x=1mm, y=1mm]
  \bkthrcgrid{-5}{5}{-4}{4}{4.6}{2}
  \bkthrcpt{(-4,2)}{4.6}{P}{south}{bkRed}
  \bkthrcpt{(1,4)}{4.6}{Q}{west}{bkBlue}
  \bkthrcpt{(3,-3)}{4.6}{R}{west}{bkGreen!80!black}
  \bkthrcpt{(-1,-4)}{4.6}{S}{east}{bkPurple}
\end{tikzpicture}
\end{minipage}
\bkjans{（1）\nobreak P$(-4,\ 2)$、Q$(1,\ 4)$、R$(3,\ -3)$、S$(-1,\ -4)$\quad（2）\nobreak \R{右}{みぎ}へ3、そこから\R{下}{した}へ2}
\end{jibun}

\begin{reidai}[どの\R{区画}{く|かく}にある？]
\R{次}{つぎ}の\R{点}{てん}は、\R{右上}{みぎ|うえ}・\R{左上}{ひだり|うえ}・\R{左下}{ひだり|した}・\R{右下}{みぎ|した}のどの\R{区画}{く|かく}にありますか。\R{軸}{じく}の\R{上}{うえ}にあるときは、どちらの\R{軸}{じく}の\R{上}{うえ}かを\R{答}{こた}えましょう。\par
（1）$(-4,\ 7)$\quad（2）$(5,\ -1)$\quad（3）$(-2,\ -6)$\quad（4）$(0,\ -3)$\par
\kangaekata $x$ \R{座標}{ざ|ひょう}の\R{符号}{ふ|ごう}で\R{左右}{さ|ゆう}、$y$ \R{座標}{ざ|ひょう}の\R{符号}{ふ|ごう}で\R{上下}{じょう|げ}が\R{決}{き}まります。0なら\R{軸}{じく}の\R{上}{うえ}です。\par
\shiki（1）\nobreak $x$ は－で\R{左}{ひだり}、$y$ は＋で\R{上}{うえ}\quad（2）\nobreak $x$ は＋で\R{右}{みぎ}、$y$ は－で\R{下}{した}\par
\hspace*{2.6\zw}（3）\nobreak \R{左}{ひだり}と\R{下}{した}\quad（4）\nobreak $x=0$ なので、\R{左右}{さ|ゆう}には\R{動}{うご}かない。\R{下}{した}へ3\par
\kotae（1）\nobreak \R{左上}{ひだり|うえ}\quad（2）\nobreak \R{右下}{みぎ|した}\quad（3）\nobreak \R{左下}{ひだり|した}\quad（4）\nobreak $y$ \R{軸}{じく}の\R{上}{うえ}
\end{reidai}

\begin{jibun}
\R{次}{つぎ}の\R{点}{てん}は、どの\R{区画}{く|かく}にありますか。\R{軸}{じく}の\R{上}{うえ}なら、どちらの\R{軸}{じく}の\R{上}{うえ}かを\R{答}{こた}えましょう。（1）$(3,\ 8)$（2）$(-2,\ -6)$（3）$(6,\ 0)$（4）$(-5,\ 4)$
\bkjans{（1）\nobreak \R{右上}{みぎ|うえ}\quad（2）\nobreak \R{左下}{ひだり|した}\quad（3）\nobreak $x$ \R{軸}{じく}の\R{上}{うえ}\quad（4）\nobreak \R{左上}{ひだり|うえ}}
\end{jibun}

\begin{reidai}[\R{点}{てん}を\R{動}{うご}かす]
（1）\R{点}{てん} A$(2,\ -1)$ を、\R{右}{みぎ}へ4、\R{上}{うえ}へ3\R{動}{うご}かした\R{点}{てん} B と、\R{左}{ひだり}へ5、\R{下}{した}へ2\R{動}{うご}かした\R{点}{てん} C の\R{座標}{ざ|ひょう}を\R{求}{もと}めましょう。\par
（2）ゲームの\R{画面}{が|めん}は、\R{原点}{げん|てん}が\R{左上}{ひだり|うえ}のすみで、$x$ は\R{右}{みぎ}\R{向}{む}き、$y$ は\R{下}{した}\R{向}{む}きが＋です。キャラクターは $(100,\ 50)$ にいて、ボタンを1\R{回}{かい}おすと\R{右}{みぎ}へ20、\R{上}{うえ}へ15\R{動}{うご}きます。3\R{回}{かい}おすと、どこへ\R{行}{い}きますか。\par
\kangaekata \R{右}{みぎ}・\R{左}{ひだり}は $x$ \R{座標}{ざ|ひょう}、\R{上}{うえ}・\R{下}{した}は $y$ \R{座標}{ざ|ひょう}だけが\R{変}{か}わります。（2）は $y$ が\R{下}{した}\R{向}{む}きに＋なので、\R{上}{うえ}へ\R{動}{うご}くと $y$ はへります。\par
\shiki（1）\nobreak B：$(2+4,\ -1+3)=(6,\ 2)$\par
\hspace*{2.6\zw}C：$(2-5,\ -1-2)=(-3,\ -3)$\par
\hspace*{2.6\zw}（2）\nobreak 3\R{回}{かい}で\R{右}{みぎ}へ $20\times3=60$、\R{上}{うえ}へ $15\times3=45$\par
\hspace*{2.6\zw}$(100+60,\ 50-45)=(160,\ 5)$\par
\kotae（1）\nobreak B$(6,\ 2)$、C$(-3,\ -3)$\quad（2）\nobreak $(160,\ 5)$
\end{reidai}

\begin{jibun}
（1）\R{点}{てん} A$(-1,\ 3)$ を、\R{右}{みぎ}へ2、\R{下}{した}へ5\R{動}{うご}かした\R{点}{てん} B と、\R{左}{ひだり}へ3、\R{上}{うえ}へ1\R{動}{うご}かした\R{点}{てん} C の\R{座標}{ざ|ひょう}を\R{求}{もと}めましょう。（2）\R{原点}{げん|てん}が\R{左上}{ひだり|うえ}、$y$ が\R{下}{した}\R{向}{む}きに＋の\R{画面}{が|めん}で、$(60,\ 80)$ にいるキャラクターが、1\R{回}{かい}で\R{左}{ひだり}へ10、\R{下}{した}へ25\R{動}{うご}きます。2\R{回}{かい}\R{動}{うご}くと、どこへ\R{行}{い}きますか。
\bkjans{（1）\nobreak B$(1,\ -2)$、C$(-4,\ 4)$\quad（2）\nobreak $(60-20,\ 80+50)=(40,\ 130)$}
\end{jibun}

\begin{reidai}[\R{下}{した}つき\R{文字}{も|じ}と、\R{動}{うご}いた\R{量}{りょう}]
ロボットの\R{手}{て}の\R{先}{さき}が、\R{点}{てん} $(x_1,\ y_1)=(3,\ -2)$ から\R{点}{てん} $(x_2,\ y_2)=(-1,\ 4)$ へ\R{動}{うご}きました（\R{単位}{たん|い}は cm）。\par
（1）$x_1$ と $y_2$ はいくつですか。（2）$x_2-x_1$ と $y_2-y_1$ を\R{求}{もと}めましょう。（3）\R{手}{て}の\R{先}{さき}は、\R{左右}{さ|ゆう}・\R{上下}{じょう|げ}にどれだけ\R{動}{うご}きましたか。\par
\kangaekata 1は1\R{番}{ばん}め（はじめ）、2は2\R{番}{ばん}め（あと）の\R{名札}{な|ふだ}です。\R{動}{うご}いた\R{量}{りょう}は「あと－まえ」。マイナスなら\R{左}{ひだり}や\R{下}{した}です。\par
\shiki（1）\nobreak $x_1$ ははじめの $x$ \R{座標}{ざ|ひょう}で3。$y_2$ はあとの $y$ \R{座標}{ざ|ひょう}で4\par
\hspace*{2.6\zw}（2）\nobreak $x_2-x_1=-1-3=-4$、$y_2-y_1=4-(-2)=4+2=6$\par
\hspace*{2.6\zw}（3）\nobreak $-4$ は\R{左}{ひだり}へ4、6は\R{上}{うえ}へ6\par
\kotae（1）\nobreak $x_1=3$、$y_2=4$\quad（2）\nobreak $-4$ と6\quad（3）\nobreak \R{左}{ひだり}へ\qty{4}{cm}、\R{上}{うえ}へ\qty{6}{cm}
\end{reidai}

\begin{jibun}
ロボットの\R{手}{て}の\R{先}{さき}が、$(x_1,\ y_1)=(-2,\ 5)$ から $(x_2,\ y_2)=(4,\ 1)$ へ\R{動}{うご}きました（cm）。（1）$y_1$ と $x_2$ はいくつですか。（2）$x_2-x_1$ と $y_2-y_1$ を\R{求}{もと}めましょう。（3）\R{左右}{さ|ゆう}・\R{上下}{じょう|げ}にどれだけ\R{動}{うご}きましたか。
\bkjans{（1）\nobreak $y_1=5$、$x_2=4$\quad（2）\nobreak $4-(-2)=6$、$1-5=-4$\quad（3）\nobreak \R{右}{みぎ}へ\qty{6}{cm}、\R{下}{した}へ\qty{4}{cm}}
\end{jibun}

\begin{reidai}[\R{原点}{げん|てん}をとりかえる]
\noindent\begin{minipage}[c]{0.49\linewidth}
\R{駅}{えき}を\R{原点}{げん|てん}、\R{東}{ひがし}を $x$ の＋、\R{北}{きた}を $y$ の＋とし、1\R{目}{め}もりを\qty{100}{m}とします。\R{学校}{がっ|こう}は $(5,\ 2)$、\R{家}{いえ}は $(-3,\ 4)$ です。\R{向}{む}きはそのままで、\R{原点}{げん|てん}を\R{家}{いえ}にとりかえると、\R{学校}{がっ|こう}と\R{駅}{えき}の\R{座標}{ざ|ひょう}はどうなりますか。
\end{minipage}\hfill
\begin{minipage}[c]{0.47\linewidth}\centering
\begin{tikzpicture}[x=1mm, y=1mm]
  \bkthrcgrid{-4}{6}{-1}{5}{5}{0}
  \bkthrcpt{(5,2)}{5}{\R{学校}{がっ|こう}}{north}{bkBlue}
  \bkthrcpt{(-3,4)}{5}{\R{家}{いえ}}{south}{bkOrange}
  \node[anchor=south west, font=\sffamily\footnotesize, fill=white, inner sep=0.6pt] at (0.5,0.5) {\R{駅}{えき}};
  \draw[bkGreen!70!black, line width=1.1pt, -{Stealth[length=4pt]}] (-15,20) -- (24,20);
  \draw[bkGreen!70!black, line width=1.1pt, -{Stealth[length=4pt]}] (25,19.5) -- (25,11.5);
\end{tikzpicture}
\end{minipage}
\par
\kangaekata \R{新}{あたら}しい\R{座標}{ざ|ひょう}は「\R{家}{いえ}から\R{見}{み}て、どれだけ\R{東}{ひがし}・\R{北}{きた}にあるか」です。\R{家}{いえ}から\R{学校}{がっ|こう}への\R{動}{うご}きを「あと－まえ」で\R{求}{もと}めます（\R{緑}{みどり}の\R{矢印}{や|じるし}）。\par
\shiki \R{学校}{がっ|こう}：$x$ は $5-(-3)=8$、$y$ は $2-4=-2$\par
\hspace*{2.6\zw}\R{駅}{えき}：$x$ は $0-(-3)=3$、$y$ は $0-4=-4$\par
\hspace*{2.6\zw}\R{確}{たし}かめ：マスを\R{数}{かぞ}えると、\R{家}{いえ}から\R{学校}{がっ|こう}は\R{東}{ひがし}へ8、\R{南}{みなみ}へ2\par
\kotae \R{学校}{がっ|こう} $(8,\ -2)$、\R{駅}{えき} $(3,\ -4)$
\end{reidai}

\begin{jibun}
\R{駅}{えき}を\R{原点}{げん|てん}、\R{東}{ひがし}を $x$ の＋、\R{北}{きた}を $y$ の＋とすると、\R{公園}{こう|えん}は $(4,\ -1)$、\R{家}{いえ}は $(-2,\ 3)$ です。\R{向}{む}きはそのままで\R{原点}{げん|てん}を\R{家}{いえ}にとりかえると、\R{公園}{こう|えん}と\R{駅}{えき}の\R{座標}{ざ|ひょう}はどうなりますか。
\bkjans{\R{公園}{こう|えん} $(4-(-2),\ -1-3)=(6,\ -4)$、\R{駅}{えき} $(0-(-2),\ 0-3)=(2,\ -3)$}
\end{jibun}

\begin{reidai}[\R{絶対値}{ぜっ|たい|ち}]
\R{次}{つぎ}の\R{値}{あたい}を\R{求}{もと}めましょう。（1）$\lvert -7\rvert $\quad（2）$\lvert 2.5\rvert $\quad（3）$\lvert 0\rvert $\quad（4）$\lvert -3\rvert +\lvert 2\rvert $\quad（5）$\lvert -6\rvert -\lvert 4\rvert $\par
\kangaekata \R{絶対値}{ぜっ|たい|ち}は0からの\R{距離}{きょ|り}なので、マイナスの\R{符号}{ふ|ごう}をとります。（4）（5）は、それぞれの\R{絶対値}{ぜっ|たい|ち}を\R{先}{さき}に\R{求}{もと}めてから\R{計算}{けい|さん}します。\par
\shiki（1）\nobreak $-7$ は0から7\quad（2）\nobreak 2.5は0から2.5\quad（3）\nobreak 0は0から0\par
\hspace*{2.6\zw}（4）\nobreak $\lvert -3\rvert +\lvert 2\rvert =3+2=5$\quad（5）\nobreak $\lvert -6\rvert -\lvert 4\rvert =6-4=2$\par
\kotae（1）\nobreak 7\quad（2）\nobreak 2.5\quad（3）\nobreak 0\quad（4）\nobreak 5\quad（5）\nobreak 2
\end{reidai}

\begin{jibun}
\R{次}{つぎ}の\R{値}{あたい}を\R{求}{もと}めましょう。（1）$\lvert -12\rvert $\quad（2）$\lvert 0.8\rvert $\quad（3）$\lvert -1.5\rvert $\quad（4）$\lvert -5\rvert +\lvert -3\rvert $\quad（5）$\lvert -9\rvert -\lvert 5\rvert $
\bkjans{（1）\nobreak 12\quad（2）\nobreak 0.8\quad（3）\nobreak 1.5\quad（4）\nobreak $5+3=8$\quad（5）\nobreak $9-5=4$}
\end{jibun}

\begin{reidai}[\R{横}{よこ}・たてに\R{並}{なら}んだ2\R{点}{てん}の\R{距離}{きょ|り}]
（1）A$(-3,\ 2)$ と B$(4,\ 2)$ の\R{距離}{きょ|り}を\R{求}{もと}めましょう。（2）C$(1,\ -5)$ と D$(1,\ 3)$ の\R{距離}{きょ|り}を\R{求}{もと}めましょう。\par
\kangaekata（1）は $y$ \R{座標}{ざ|ひょう}が\R{同}{おな}じなので\R{横}{よこ}に\R{並}{なら}び、$x$ \R{座標}{ざ|ひょう}の\R{差}{さ}の\R{絶対値}{ぜっ|たい|ち}。（2）は $x$ \R{座標}{ざ|ひょう}が\R{同}{おな}じなので、たてに\R{並}{なら}び、$y$ \R{座標}{ざ|ひょう}の\R{差}{さ}の\R{絶対値}{ぜっ|たい|ち}です。\par
\shiki（1）\nobreak $\lvert 4-(-3)\rvert =\lvert 4+3\rvert =\lvert 7\rvert =7$\par
\hspace*{2.6\zw}（2）\nobreak $\lvert 3-(-5)\rvert =\lvert 3+5\rvert =\lvert 8\rvert =8$\par
\hspace*{2.6\zw}\R{確}{たし}かめ：\R{方眼}{ほう|がん}にかいてマスを\R{数}{かぞ}えても、7と8\par
\kotae（1）\nobreak 7\quad（2）\nobreak 8
\end{reidai}

\begin{jibun}
（1）A$(-6,\ -1)$ と B$(-2,\ -1)$ の\R{距離}{きょ|り}を\R{求}{もと}めましょう。（2）C$(2,\ 4)$ と D$(2,\ -3)$ の\R{距離}{きょ|り}を\R{求}{もと}めましょう。
\bkjans{（1）\nobreak $\lvert -2-(-6)\rvert =\lvert 4\rvert =4$\quad（2）\nobreak $\lvert -3-4\rvert =\lvert -7\rvert =7$}
\end{jibun}

\begin{reidai}[\R{長方形}{ちょう|ほう|けい}の\R{頂点}{ちょう|てん}・まわりの\R{長}{なが}さ・\R{面積}{めん|せき}]
\noindent\begin{minipage}[c]{0.54\linewidth}
\R{長方形}{ちょう|ほう|けい} ABCD の3つの\R{頂点}{ちょう|てん}が A$(-3,\ -2)$、B$(5,\ -2)$、C$(5,\ 4)$ です（1\R{目}{め}もり\qty{1}{cm}）。（1）\R{頂点}{ちょう|てん} D の\R{座標}{ざ|ひょう}を\R{求}{もと}めましょう。（2）\R{長方形}{ちょう|ほう|けい}のまわりの\R{長}{なが}さと\R{面積}{めん|せき}を\R{求}{もと}めましょう。
\end{minipage}\hfill
\begin{minipage}[c]{0.42\linewidth}\centering
\begin{tikzpicture}[x=1mm, y=1mm]
  \bkthrcgrid{-4}{6}{-3}{5}{4.6}{0}
  \fill[bkBlue!15] (-13.8,-9.2) rectangle (23,18.4);
  \draw[bkBlue, line width=1pt] (-13.8,-9.2) rectangle (23,18.4);
  \bkthrcpt{(-3,-2)}{4.6}{A}{north east}{bkBlue}
  \bkthrcpt{(5,-2)}{4.6}{B}{north west}{bkBlue}
  \bkthrcpt{(5,4)}{4.6}{C}{south west}{bkBlue}
  \bkthrcpt{(-3,4)}{4.6}{D}{south east}{bkRed}
\end{tikzpicture}
\end{minipage}
\par
\kangaekata D は、A と\R{同}{おな}じたての\R{線}{せん}の\R{上}{うえ}（$x$ が\R{同}{おな}じ）、C と\R{同}{おな}じ\R{横}{よこ}の\R{線}{せん}の\R{上}{うえ}（$y$ が\R{同}{おな}じ）にあります。\R{辺}{へん}の\R{長}{なが}さは、\R{横}{よこ}・たてに\R{並}{なら}んだ2\R{点}{てん}の\R{距離}{きょ|り}です。\par
\shiki（1）\nobreak D の $x$ は A と\R{同}{おな}じ $-3$、$y$ は C と\R{同}{おな}じ4\par
\hspace*{2.6\zw}（2）\nobreak \R{横}{よこ}の\R{辺}{へん} AB は $\lvert 5-(-3)\rvert =8$、たての\R{辺}{へん} BC は $\lvert 4-(-2)\rvert =6$\par
\hspace*{2.6\zw}まわり：$(8+6)\times2=28$、\R{面積}{めん|せき}：$8\times6=48$\par
\kotae（1）\nobreak D$(-3,\ 4)$\quad（2）\nobreak まわり\qty{28}{cm}、\R{面積}{めん|せき}\qty{48}{cm^2}
\end{reidai}

\begin{jibun}
\R{長方形}{ちょう|ほう|けい} ABCD の3つの\R{頂点}{ちょう|てん}が A$(-1,\ -3)$、B$(3,\ -3)$、C$(3,\ 2)$ です（1\R{目}{め}もり\qty{1}{cm}）。（1）D の\R{座標}{ざ|ひょう}を\R{求}{もと}めましょう。（2）まわりの\R{長}{なが}さと\R{面積}{めん|せき}を\R{求}{もと}めましょう。
\bkjans{（1）\nobreak D$(-1,\ 2)$\quad（2）\nobreak AB は4、BC は5。まわり $(4+5)\times2=18$ で\qty{18}{cm}、\R{面積}{めん|せき} $4\times5=20$ で\qty{20}{cm^2}}
\end{jibun}

%% ---------------- まちがえやすいところ ----------------
\bksection{tsumazuki}{まちがえやすいところをチェック}\label{bk:3-2:tsuma}

\begin{tsumazuki}{$x$ と $y$ の\R{順番}{じゅん|ばん}}
\bkng{\R{点}{てん} $(3,\ 2)$ を、\R{上}{うえ}へ3、\R{右}{みぎ}へ2\R{進}{すす}んでかく}
\bkok{かっこの\R{前}{まえ}が $x$（\R{右}{みぎ}へ3）、\R{後}{あと}が $y$（\R{上}{うえ}へ2）}
\R{入}{い}れかえると、$(2,\ 3)$ というちがう\R{点}{てん}になります。「\R{横}{よこ}が\R{先}{さき}、たてが\R{後}{あと}」と\R{覚}{おぼ}えよう。
\end{tsumazuki}

\begin{tsumazuki}{マイナスの\R{座標}{ざ|ひょう}の\R{向}{む}き}
\bkng{\R{点}{てん} $(-3,\ 2)$ を、\R{右}{みぎ}へ3、\R{上}{うえ}へ2\R{進}{すす}んでかく}
\bkok{$x$ がマイナスなので、\R{左}{ひだり}へ3、そこから\R{上}{うえ}へ2}
マイナスは「\R{反対}{はん|たい}の\R{向}{む}き」（2-2）。$x$ なら\R{左}{ひだり}、$y$ なら\R{下}{した}です。
\end{tsumazuki}

\begin{tsumazuki}{\R{下}{した}つき\R{文字}{も|じ}をかけ\R{算}{ざん}にする}
\bkng{$x_1=x\times1=x$}
\bkok{$x_1$ は「1\R{番}{ばん}めの $x$」。$x_1$ と $x_2$ は、ちがう\R{数}{かず}でもよい}
\R{右下}{みぎ|した}の\R{小}{ちい}さな\R{数字}{すう|じ}は\R{名札}{な|ふだ}です。\R{右上}{みぎ|うえ}の $x^2$（2\R{乗}{じょう}）とも、\R{前}{まえ}の $2x$ ともちがいます。
\end{tsumazuki}

\begin{tsumazuki}{マイナスをはさむ2\R{点}{てん}の\R{距離}{きょ|り}}
\bkng{A$(-2,\ 1)$ と B$(3,\ 1)$ の\R{距離}{きょ|り}は、$3-2=1$}
\bkok{$3-(-2)=3+2=5$。マスを\R{数}{かぞ}えても5}
$-2$ のマイナスを\R{落}{お}としてはいけません。0をまたぐときは、0までの\R{距離}{きょ|り}2と、0からの\R{距離}{きょ|り}3をたした5になります。
\end{tsumazuki}

\begin{tsumazuki}{\R{絶対値}{ぜっ|たい|ち}の\R{中}{なか}に\R{式}{しき}があるとき}
\bkng{$\lvert -3+1\rvert =3+1=4$}
\bkok{\R{中}{なか}を\R{先}{さき}に\R{計算}{けい|さん}して、$\lvert -3+1\rvert =\lvert -2\rvert =2$}
たての\R{線}{せん}ではさまれた\R{所}{ところ}は、かっこと\R{同}{おな}じ「ひとまとまり」です（2-11）。
\end{tsumazuki}

\begin{tsumazuki}{\R{途中}{と|ちゅう}で\R{原点}{げん|てん}を\R{変}{か}える}
\bkng{\R{手}{て}は\R{地面}{じ|めん}から1.5、\R{最高点}{さい|こう|てん}は\R{手}{て}から5.0。\R{上}{あ}がった\R{高}{たか}さは $5.0-1.5=3.5$}
\bkok{\R{地面}{じ|めん}を\R{原点}{げん|てん}にそろえて、$6.5-1.5=5.0$（\R{手}{て}を\R{原点}{げん|てん}にそろえても $5.0-0=5.0$）}
1つの\R{問題}{もん|だい}の\R{中}{なか}では、\R{原点}{げん|てん}と＋の\R{向}{む}きを1つに\R{決}{き}め、\R{図}{ず}にかいておこう。
\end{tsumazuki}

%% ---------------- 物理・工学 ----------------
\begin{butsuri}
\noindent\begin{minipage}[c]{0.53\linewidth}
\par\addvspace{3pt}\textbf{ロボットの\R{手}{て}の\R{先}{さき}の\R{位置}{い|ち}}\par\noindent
\R{工場}{こう|じょう}で\R{部品}{ぶ|ひん}を\R{運}{はこ}ぶロボットは、\R{台}{だい}の\R{中心}{ちゅう|しん}を\R{原点}{げん|てん}にして、\R{手}{て}の\R{先}{さき}の\R{位置}{い|ち}を $(x,\ y)$（\R{単位}{たん|い}は mm）で\R{覚}{おぼ}えています。「$(250,\ 120)$ の\R{部品}{ぶ|ひん}をつかんで、$(-180,\ 300)$ へ\R{運}{はこ}べ」という\R{命令}{めい|れい}を\R{受}{う}けると、$x$ を $-180-250=-430$、$y$ を $300-120=180$ だけ\R{動}{うご}かせばよいと\R{計算}{けい|さん}します。
\end{minipage}\hfill
\begin{minipage}[c]{0.43\linewidth}\centering
\begin{tikzpicture}[x=1.2mm, y=1.2mm]
  \draw[-{Stealth[length=3pt]}, bkInk, line width=0.6pt] (-4,0) -- (38,0) node[anchor=west, font=\small, inner sep=1pt] {$x$};
  \draw[-{Stealth[length=3pt]}, bkInk, line width=0.6pt] (0,-4) -- (0,30) node[anchor=south, font=\small, inner sep=1pt] {$y$};
  \fill[bkGray!70] (-3,-3) rectangle (3,1);
  \draw[bkOrange!80!black, line width=3.4pt, line cap=round] (0,0) -- (14,17) -- (29,10);
  \fill[bkGray] (0,0) circle (1.8) (14,17) circle (1.5);
  \draw[bkInk, line width=0.9pt] (29,10) -- (32,12) (29,10) -- (32,8);
  \fill[bkRed] (31,10) circle (0.9);
  \draw[bkRed, line width=0.5pt, densely dashed] (31,10) -- (31,0) (31,10) -- (0,10);
  \node[anchor=north, font=\sffamily\footnotesize, inner sep=1pt] at (2,-1) {O};
  \node[anchor=south, font=\sffamily\footnotesize, text=bkRed!80!black, inner sep=1pt] at (31,13) {\R{手}{て}の\R{先}{さき}};
\end{tikzpicture}
\end{minipage}

\par\addvspace{3pt}\textbf{\R{画面}{が|めん}の\R{座標}{ざ|ひょう}は、$y$ が\R{下}{した}\R{向}{む}き}\label{bk:3-2:gamen}\par\noindent
\begin{minipage}[c]{0.56\linewidth}
スマホやパソコンの\R{画面}{が|めん}は、ピクセルという\R{小}{ちい}さな\R{点}{てん}が\R{並}{なら}んでできています。\R{横}{よこ}1920\R{個}{こ}・たて1080\R{個}{こ}の\R{画面}{が|めん}なら、\R{左上}{ひだり|うえ}のすみが\R{原点}{げん|てん} $(0,\ 0)$ で、$x$ は\R{右}{みぎ}へ、$y$ は\R{下}{した}へ\R{数}{かぞ}えます。\R{画面}{が|めん}は\R{上}{うえ}の\R{行}{ぎょう}から\R{順}{じゅん}に\R{表示}{ひょう|じ}するので、\R{上}{うえ}から\R{数}{かぞ}えるほうが\R{便利}{べん|り}だからです。\R{右下}{みぎ|した}のすみは $(1919,\ 1079)$ です（0から\R{数}{かぞ}えるので、1920ではなく1919）。
\end{minipage}\hfill
\begin{minipage}[c]{0.4\linewidth}\centering
\begin{tikzpicture}[x=1.3mm, y=1.3mm]
  \filldraw[fill=bkInk!85, draw=bkInk, line width=0.8pt, rounded corners=1mm] (-2,-24) rectangle (38,2);
  \fill[bkBlue!25] (0,-22) rectangle (36,0);
  \draw[-{Stealth[length=3pt]}, bkRed, line width=0.8pt] (0.5,-0.5) -- (14,-0.5);
  \draw[-{Stealth[length=3pt]}, bkRed, line width=0.8pt] (0.5,-0.5) -- (0.5,-13);
  \node[anchor=north west, font=\small, text=bkRed!80!black, inner sep=1pt] at (14,-1) {$x$};
  \node[anchor=north west, font=\small, text=bkRed!80!black, inner sep=1pt] at (1.5,-12) {$y$};
  \node[anchor=south, font=\sffamily\footnotesize, inner sep=1pt] at (2,2.5) {$(0,\ 0)$};
  \fill[bkOrange] (24,-14) circle (1.4);
  \node[anchor=south, font=\sffamily\footnotesize, inner sep=1pt, fill=bkBlue!25] at (24,-12.5) {$(1200,\ 700)$};
\end{tikzpicture}
\end{minipage}

\end{butsuri}
\begin{butsuri}
\par\addvspace{3pt}\textbf{\R{工作機械}{こう|さく|き|かい}の\R{原点}{げん|てん}}\par\noindent
\begin{minipage}[c]{0.5\linewidth}
\R{金属}{きん|ぞく}の\R{板}{いた}に\R{穴}{あな}をあける\R{機械}{き|かい}（NC フライス\R{盤}{ばん}）では、\R{板}{いた}の\R{左下}{ひだり|した}の\R{角}{かど}を\R{原点}{げん|てん}にして、「$(20,\ 15)$ と $(60,\ 15)$ に\R{穴}{あな}をあける」のように\R{座標}{ざ|ひょう}（mm）で\R{命令}{めい|れい}します。\R{角}{かど}を\R{原点}{げん|てん}にすると、\R{図面}{ず|めん}に\R{書}{か}かれた「\R{角}{かど}から\R{何}{なん}mm」がそのまま\R{座標}{ざ|ひょう}になり、\R{計算}{けい|さん}がいりません。\R{原点}{げん|てん}を\R{計算}{けい|さん}が\R{楽}{らく}な\R{所}{ところ}に\R{置}{お}く、よい\R{例}{れい}です。
\end{minipage}\hfill
\begin{minipage}[c]{0.47\linewidth}\centering
\begin{tikzpicture}[x=0.5mm, y=0.5mm]
  \filldraw[fill=bkGray!25, draw=bkGray!90!black, line width=0.8pt] (0,0) rectangle (80,40);
  \foreach \hx in {20,60} {\filldraw[fill=white, draw=bkInk, line width=0.7pt] (\hx,15) circle (4.5);
    \draw[bkRed, line width=0.4pt] (\hx-6,15) -- (\hx+6,15) (\hx,9) -- (\hx,21);}
  \draw[-{Stealth[length=3pt]}, bkInk, line width=0.7pt] (0,0) -- (92,0);
  \draw[-{Stealth[length=3pt]}, bkInk, line width=0.7pt] (0,0) -- (0,52);
  \node[anchor=west, font=\small, inner sep=1pt] at (92,0) {$x$};
  \node[anchor=south, font=\small, inner sep=1pt] at (0,52) {$y$};
  \node[anchor=north east, font=\sffamily\footnotesize, inner sep=1pt] at (-1,-1) {O};
  \draw[{Stealth[length=2.5pt]}-{Stealth[length=2.5pt]}, bkBlue, line width=0.5pt] (0,-8) -- (20,-8);
  \node[anchor=north, font=\sffamily\footnotesize, text=bkBlue!80!black, inner sep=1pt] at (10,-9) {20};
  \draw[{Stealth[length=2.5pt]}-{Stealth[length=2.5pt]}, bkBlue, line width=0.5pt] (0,-20) -- (60,-20);
  \node[anchor=north, font=\sffamily\footnotesize, text=bkBlue!80!black, inner sep=1pt] at (30,-21) {60};
  \draw[{Stealth[length=2.5pt]}-{Stealth[length=2.5pt]}, bkBlue, line width=0.5pt] (-8,0) -- (-8,15);
  \node[anchor=east, font=\sffamily\footnotesize, text=bkBlue!80!black, inner sep=1pt] at (-9,7.5) {15};
  \node[anchor=south, font=\sffamily\footnotesize, inner sep=1pt] at (40,42) {\R{金属}{きん|ぞく}の\R{板}{いた}（mm）};
\end{tikzpicture}
\end{minipage}

\par\addvspace{3pt}\textbf{ロボットそうじ\R{機}{き}の\R{地図}{ち|ず}}\par\noindent
\begin{minipage}[c]{0.56\linewidth}
ロボットそうじ\R{機}{き}の\R{中}{なか}には、\R{充電台}{じゅう|でん|だい}を\R{原点}{げん|てん}にして、\R{床}{ゆか}を\R{小}{ちい}さな\R{正方形}{せい|ほう|けい}のマスに\R{分}{わ}けた\R{地図}{ち|ず}を\R{作}{つく}るものがあります。\R{走}{はし}りながら、\R{自分}{じ|ぶん}がいるマスの\R{座標}{ざ|ひょう}と、\R{壁}{かべ}や\R{家具}{か|ぐ}にぶつかったマスの\R{座標}{ざ|ひょう}を\R{記録}{き|ろく}していきます。まだ\R{通}{とお}っていないマスの\R{座標}{ざ|ひょう}がわかるので、そうじのしのこしをへらせます。
\end{minipage}\hfill
\begin{minipage}[c]{0.4\linewidth}\centering
\begin{tikzpicture}[x=1mm, y=1mm]
  \foreach \cx/\cy in {0/0,1/0,2/0,3/0,4/0,5/0,5/1,4/1,3/1,2/1,1/1,0/1,0/2,1/2,2/2} {\fill[bkGreen!25] (\cx*6,\cy*6) rectangle ++(6,6);}
  \foreach \cx/\cy in {3/2,4/2,3/3,4/3} {\fill[bkGray!60] (\cx*6,\cy*6) rectangle ++(6,6);}
  \draw[bkLine, line width=0.4pt] (0,0) grid[step=6] (42,30);
  \draw[bkInk, line width=1pt] (0,0) rectangle (42,30);
  \draw[bkOrange!80!black, line width=0.9pt, -{Stealth[length=3pt]}] (3,3) -- (33,3) -- (33,9) -- (3,9) -- (3,15) -- (15,15);
  \fill[bkNavy] (0,0) circle (1.4);
  \node[anchor=north, font=\sffamily\footnotesize, inner sep=1pt] at (2,-1) {\R{充電台}{じゅう|でん|だい}（\R{原点}{げん|てん}）};
  \node[font=\sffamily\footnotesize, fill=bkGray!60, inner sep=0.6pt] at (24,21) {つくえ};
\end{tikzpicture}
\end{minipage}

\par\addvspace{3pt}\textbf{\R{大学}{だい|がく}への\R{種}{たね}}\par\noindent
\R{空間}{くう|かん}の\R{中}{なか}の\R{場所}{ば|しょ}は、\R{高}{たか}さの $z$ を\R{加}{くわ}えた3つの\R{数}{かず} $(x,\ y,\ z)$ で\R{表}{あらわ}します（4-13）。また、「\R{原点}{げん|てん}からの\R{距離}{きょ|り}と\R{向}{む}きの\R{角度}{かく|ど}」の2つで\R{表}{あらわ}す\R{方法}{ほう|ほう}もあります（\R{極座標}{きょく|ざ|ひょう}、4-2）。どれも、この\R{章}{しょう}の「\R{原点}{げん|てん}と\R{向}{む}きを\R{決}{き}めて、\R{数}{かず}で\R{場所}{ば|しょ}を\R{言}{い}う」\R{考}{かんが}え\R{方}{かた}の\R{続}{つづ}きです。
\end{butsuri}

\begin{motto}[2]
\textbf{\R{部屋}{へ|や}の\R{中}{なか}を\R{飛}{と}ぶハエは、3つの\R{数}{かず}}\par
\noindent\begin{minipage}[c]{0.6\linewidth}
\R{天井}{てん|じょう}のハエは2つの\R{数}{かず}で\R{言}{い}えました。\R{部屋}{へ|や}の\R{中}{なか}を\R{飛}{と}ぶハエは、\R{床}{ゆか}のすみを\R{原点}{げん|てん}にして、\R{床}{ゆか}にそった $x$ と $y$ の\R{向}{む}きに、\R{上}{うえ}\R{向}{む}きの $z$ を\R{加}{くわ}えた3つの\R{数}{かず} $(x,\ y,\ z)$ で\R{場所}{ば|しょ}が\R{決}{き}まります。$(3,\ 2,\ 4)$ は、$x$ の\R{向}{む}きに3、$y$ の\R{向}{む}きに2\R{進}{すす}み、そこから\R{上}{うえ}へ4の\R{場所}{ば|しょ}です（くわしくは 4-13）。
\end{minipage}\hfill
\begin{minipage}[c]{0.37\linewidth}\centering
\begin{tikzpicture}[x={(0.9mm,0mm)}, y={(0.5mm,0.4mm)}, z={(0mm,1mm)}]
  \draw[bkLine, line width=0.4pt] (0,0,0) -- (0,24,0);
  \foreach \k in {6,12,...,36} \draw[bkLine, line width=0.4pt] (\k,0,0) -- (\k,24,0);
  \foreach \k in {6,12,18,24} \draw[bkLine, line width=0.4pt] (0,\k,0) -- (36,\k,0);
  \draw[-{Stealth[length=3pt]}, bkInk, line width=0.7pt] (0,0,0) -- (40,0,0) node[anchor=west, font=\small, inner sep=1pt] {$x$};
  \draw[-{Stealth[length=3pt]}, bkInk, line width=0.7pt] (0,0,0) -- (0,28,0) node[anchor=south west, font=\small, inner sep=1pt] {$y$};
  \draw[-{Stealth[length=3pt]}, bkInk, line width=0.7pt] (0,0,0) -- (0,0,27) node[anchor=south, font=\small, inner sep=1pt] {$z$};
  \draw[bkRed, line width=0.6pt, densely dashed] (18,0,0) -- (18,12,0) -- (0,12,0) (18,12,0) -- (18,12,20);
  \node[circle, fill=bkRed, inner sep=0pt, minimum size=1.4mm] at (18,12,0) {};
  \node[circle, fill=bkInk, inner sep=0pt, minimum size=2.6mm] at (18,12,20) {};
  \node[anchor=west, font=\sffamily\footnotesize, inner sep=1.5pt] at (18,12,20) {$(3,\ 2,\ 4)$};
\end{tikzpicture}
\end{minipage}
\end{motto}

\begin{column}{\R{座標}{ざ|ひょう}を\R{考}{かんが}えたデカルト}
\noindent\begin{minipage}[c]{0.58\linewidth}
フランスのデカルト（1596〜1650\R{年}{ねん}）は、2-11 で、わかっている\R{数}{かず}を $a$、$b$、$c$、わからない\R{数}{かず}を $x$、$y$、$z$ と\R{書}{か}いた\R{人}{ひと}です。ベッドから\R{天井}{てん|じょう}のハエを\R{見}{み}ていて、「すみから、こちらの\R{壁}{かべ}にそって\R{何}{なん}マス、あちらの\R{壁}{かべ}にそって\R{何}{なん}マス」と、ハエの\R{場所}{ば|しょ}を2つの\R{数}{かず}で\R{言}{い}えると\R{気}{き}づいた、という\R{話}{はなし}が\R{伝}{つた}わっています（ほんとうの\R{話}{はなし}かどうかは、わかりません）。
\end{minipage}\hfill
\begin{minipage}[c]{0.38\linewidth}\centering
\begin{tikzpicture}[x=1mm, y=1mm]
  \fill[bkYellowBg] (0,0) rectangle (36,24);
  \draw[bkLine, line width=0.4pt] (0,0) grid[step=6] (36,24);
  \draw[bkInk, line width=1.4pt] (0,24) -- (0,0) -- (36,0);
  \node[anchor=north east, font=\sffamily\footnotesize, inner sep=1pt] at (-0.5,-0.5) {すみ};
  \fill[bkInk] (18,12) ellipse (1.6 and 1);
  \fill[bkBlue!40] (17,13) ellipse (1.2 and 0.6) (19,13) ellipse (1.2 and 0.6);
  \draw[bkRed, line width=0.6pt, densely dashed] (0,12) -- (16.2,12) (18,0) -- (18,10.8);
  \node[anchor=north, font=\sffamily\footnotesize, text=bkRed!80!black, inner sep=1pt] at (18,-1) {3マス};
  \node[anchor=east, font=\sffamily\footnotesize, text=bkRed!80!black, inner sep=1pt] at (-1,12) {2マス};
\end{tikzpicture}
\end{minipage}
\par\vspace{2pt}
デカルトは1637\R{年}{ねん}の\R{本}{ほん}で、\R{図形}{ず|けい}を\R{数}{かず}と\R{式}{しき}で\R{表}{あらわ}して\R{計算}{けい|さん}する\R{方法}{ほう|ほう}を\R{示}{しめ}しました。それで、この\R{座標}{ざ|ひょう}は「デカルト\R{座標}{ざ|ひょう}」ともよばれます。この\R{考}{かんが}えから、3-3 で\R{学}{まな}ぶ\R{関数}{かん|すう}のグラフや、\R{第}{だい}5\R{巻}{かん}の\R{微分}{び|ぶん}・\R{積分}{せき|ぶん}が\R{生}{う}まれました。
\end{column}


%% ---------------- 章末問題 ----------------
\bkendproblems
\bkproblemgroup{\R{基本}{き|ほん}}{1}

\bkproblem \noindent\begin{minipage}[t]{0.4\linewidth}
\R{右}{みぎ}の\R{図}{ず}の\R{点}{てん} A、B、C、D の\R{座標}{ざ|ひょう}を\R{答}{こた}えましょう。
\end{minipage}\hfill
\begin{minipage}[t]{0.50\linewidth}\vspace{0mm}\centering
\begin{tikzpicture}[x=1mm, y=1mm]
  \bkthrcgrid{-5}{5}{-4}{4}{5}{2}
  \bkthrcpt{(2,4)}{5}{A}{west}{bkRed}
  \bkthrcpt{(-5,1)}{5}{B}{south}{bkBlue}
  \bkthrcpt{(-1,-3)}{5}{C}{east}{bkGreen!80!black}
  \bkthrcpt{(4,-2)}{5}{D}{west}{bkPurple}
\end{tikzpicture}
\end{minipage}
\bkpans{A$(2,\ 4)$、B$(-5,\ 1)$、C$(-1,\ -3)$、D$(4,\ -2)$}

\bkproblem \R{次}{つぎ}の\R{点}{てん}は、どの\R{区画}{く|かく}にありますか。\R{軸}{じく}の\R{上}{うえ}なら、どちらの\R{軸}{じく}の\R{上}{うえ}かを\R{答}{こた}えましょう。（1）$(-1,\ -9)$（2）$(7,\ 2)$（3）$(0,\ 5)$（4）$(4,\ -6)$
\bkpans{（1）\nobreak \R{左下}{ひだり|した}\quad（2）\nobreak \R{右上}{みぎ|うえ}\quad（3）\nobreak $y$ \R{軸}{じく}の\R{上}{うえ}\quad（4）\nobreak \R{右下}{みぎ|した}}

\bkproblem \R{点}{てん} A$(-4,\ 2)$ を、\R{右}{みぎ}へ6、\R{下}{した}へ3\R{動}{うご}かした\R{点}{てん} B と、\R{左}{ひだり}へ1、\R{上}{うえ}へ4\R{動}{うご}かした\R{点}{てん} C の\R{座標}{ざ|ひょう}を\R{求}{もと}めましょう。
\bkpans{B$(2,\ -1)$、C$(-5,\ 6)$}

\bkproblem \R{点}{てん} $(x_1,\ y_1)=(5,\ 1)$ から\R{点}{てん} $(x_2,\ y_2)=(-2,\ -3)$ へ\R{動}{うご}きました。（1）$x_2$ と $y_1$ はいくつですか。（2）$x_2-x_1$ と $y_2-y_1$ を\R{求}{もと}め、\R{左右}{さ|ゆう}・\R{上下}{じょう|げ}にどれだけ\R{動}{うご}いたかを\R{答}{こた}えましょう。
\bkpans{（1）\nobreak $x_2=-2$、$y_1=1$\quad（2）\nobreak $-2-5=-7$、$-3-1=-4$。\R{左}{ひだり}へ7、\R{下}{した}へ4}

\bkproblem \R{次}{つぎ}の\R{値}{あたい}を\R{求}{もと}めましょう。（1）$\lvert -20\rvert $\quad（2）$\lvert 3.6\rvert $\quad（3）$\lvert -4\rvert +\lvert -1\rvert $\quad（4）$\lvert -8\rvert -\lvert 6\rvert $
\bkpans{（1）\nobreak 20\quad（2）\nobreak 3.6\quad（3）\nobreak 5\quad（4）\nobreak 2}

\bkproblem（1）A$(-5,\ 3)$ と B$(2,\ 3)$ の\R{距離}{きょ|り}、（2）C$(-4,\ -6)$ と D$(-4,\ 1)$ の\R{距離}{きょ|り}を\R{求}{もと}めましょう。
\bkpans{（1）\nobreak $\lvert 2-(-5)\rvert =7$\quad（2）\nobreak $\lvert 1-(-6)\rvert =7$}

\bkproblem \R{駅}{えき}を\R{原点}{げん|てん}、\R{東}{ひがし}を $x$ の＋、\R{北}{きた}を $y$ の＋とすると、\R{図書館}{と|しょ|かん}は $(3,\ 5)$、\R{家}{いえ}は $(-4,\ 1)$ です。\R{向}{む}きはそのままで\R{原点}{げん|てん}を\R{家}{いえ}にとりかえると、\R{図書館}{と|しょ|かん}と\R{駅}{えき}の\R{座標}{ざ|ひょう}はどうなりますか。
\bkpans{\R{図書館}{と|しょ|かん} $(3-(-4),\ 5-1)=(7,\ 4)$、\R{駅}{えき} $(0-(-4),\ 0-1)=(4,\ -1)$}

\bkproblemgroup{チャレンジ}{2}

\bkproblem \R{長方形}{ちょう|ほう|けい} ABCD の3つの\R{頂点}{ちょう|てん}が A$(-4,\ -1)$、B$(2,\ -1)$、C$(2,\ 3)$ です（1\R{目}{め}もり\qty{1}{cm}）。D の\R{座標}{ざ|ひょう}と、まわりの\R{長}{なが}さ、\R{面積}{めん|せき}を\R{求}{もと}めましょう。
\bkpans{D$(-4,\ 3)$。AB $=6$、BC $=4$。まわり\qty{20}{cm}、\R{面積}{めん|せき}\qty{24}{cm^2}}

\bkproblem \R{原点}{げん|てん}が\R{左上}{ひだり|うえ}、$x$ が\R{右}{みぎ}\R{向}{む}き、$y$ が\R{下}{した}\R{向}{む}きに＋の\R{画面}{が|めん}で、ボールが $(300,\ 200)$ にあります。ボールは1\R{秒}{びょう}ごとに\R{右}{みぎ}へ40、\R{上}{うえ}へ25\R{動}{うご}きます。4\R{秒後}{びょう|ご}のボールの\R{座標}{ざ|ひょう}を\R{求}{もと}めましょう。
\bkpans{\R{右}{みぎ}へ $40\times4=160$、\R{上}{うえ}へ $25\times4=100$。$(300+160,\ 200-100)=(460,\ 100)$}

\begin{furikaeri}
\begin{checklist}
\item \R{原点}{げん|てん}と\R{軸}{じく}を\R{決}{き}めて、\R{点}{てん}をかいたり\R{座標}{ざ|ひょう}を\R{読}{よ}んだりできる。\bkfurref{bk:3-2:zahyo}
\item マイナスの\R{座標}{ざ|ひょう}、4つの\R{区画}{く|かく}、\R{点}{てん}を\R{動}{うご}かした\R{座標}{ざ|ひょう}がわかる。\bkfurref{bk:3-2:ugokasu}
\item \R{絶対値}{ぜっ|たい|ち}を\R{使}{つか}って、2\R{点}{てん}の\R{距離}{きょ|り}を\R{求}{もと}められる。\bkfurref{bk:3-2:kyori}
\end{checklist}
\end{furikaeri}

\bkfukidashi{sora}{2つの\R{数}{かず}で、\R{場所}{ば|しょ}がぴったり\R{決}{き}まるんだね。\R{次}{つぎ}の\R{章}{しょう}では、この\R{点}{てん}を\R{並}{なら}べて、\R{関数}{かん|すう}のグラフにするよ。}
